\documentclass[12pt,a4paper]{article}

\usepackage[a4paper,margin=1in]{geometry}
\usepackage{amsmath,amssymb,bm,mathtools}
\usepackage{physics}
\usepackage{siunitx}
\usepackage{setspace}
\usepackage{microtype}
\usepackage{titlesec}
\usepackage{float}
\usepackage{graphicx}
\setstretch{1.15}
\setlength{\parindent}{0pt}
\setlength{\parskip}{6pt}

\titleformat{\section}{\large\bfseries}{\thesection}{0.7em}{}
\titleformat{\subsection}{\normalsize\bfseries}{\thesubsection}{0.7em}{}

\newcommand{\ketg}{\left|g\right\rangle}
\newcommand{\kete}{\left|e\right\rangle}
\newcommand{\Efield}{\mathbf E}
\newcommand{\rvec}{\mathbf r}

\begin{document}

\begin{center}
    {\Large\bfseries Channel Estimation for a Rydberg Atomic Receiver}\\[4pt]
\end{center}

\hrule
\vspace{8pt}


% ============================================================
% PAGE 1
% ============================================================

\section{Quantum System and Two-Level Rydberg State}

A quantum system is described by a state vector. For the two-level
Rydberg-atom model, the ground state and excited state are denoted
by
\[
\ketg,
\qquad
\kete.
\]

The state of atom \(i\) is written as
\begin{equation}
\ket{\Phi_i(t)}
=
\zeta_{g,i}(t)\ketg
+
\zeta_{e,i}(t)\kete .
\tag{1}
\end{equation}

The quantities
\(\zeta_{g,i}(t)\) and \(\zeta_{e,i}(t)\) are the probability
amplitudes corresponding to the ground and excited states,
respectively.

The probability of finding the atom in the ground state is
\[
P_{g,i}(t)
=
\left|\zeta_{g,i}(t)\right|^2,
\]
while the probability of finding the atom in the excited state is
\[
P_{e,i}(t)
=
\left|\zeta_{e,i}(t)\right|^2.
\]

Since the atom must be in either the ground state or the excited
state, the total probability is unity:
\[
P_{g,i}(t)+P_{e,i}(t)=1.
\]

Therefore,
\[
\boxed{
\left|\zeta_{g,i}(t)\right|^2
+
\left|\zeta_{e,i}(t)\right|^2
=1.
}
\]

The two-level basis is chosen as
\[
\kete
=
\begin{bmatrix}
1\\
0
\end{bmatrix},
\qquad
\ketg
=
\begin{bmatrix}
0\\
1
\end{bmatrix}.
\]

Substituting these basis vectors into (1),
\[
\begin{aligned}
\ket{\Phi_i(t)}
&=
\zeta_{g,i}(t)
\begin{bmatrix}
0\\
1
\end{bmatrix}
+
\zeta_{e,i}(t)
\begin{bmatrix}
1\\
0
\end{bmatrix}\\
&=
\begin{bmatrix}
\zeta_{e,i}(t)\\
\zeta_{g,i}(t)
\end{bmatrix}.
\end{aligned}
\]

Hence, for the basis ordering
\[
(\kete,\ketg),
\]
the state vector is
\[
\ket{\Phi_i(t)}
=
\begin{bmatrix}
\zeta_{e,i}(t)\\
\zeta_{g,i}(t)
\end{bmatrix}.
\]


% ============================================================
% PAGE 2
% ============================================================

\section{Time-Dependent Schr\"odinger Equation}

The time evolution of a quantum state is described by the
time-dependent Schr\"odinger equation.

In the absence of an external interaction,
\[
i\hbar
\frac{\partial}{\partial t}
\ket{\Phi_i(t)}
=
\hat H_i
\ket{\Phi_i(t)},
\]
where \(\hat H_i\) is the free Hamiltonian of the atom. Throughout this
document, both \(i\) (in \(i\hbar\)) and \(j\) denote the imaginary unit,
\(i=j=\sqrt{-1}\); the letter \(i\) used as a subscript is the vapor-cell index.

When an external electromagnetic field interacts with the atom,
an interaction Hamiltonian \(\hat V_i\) is introduced. Therefore,
the total Hamiltonian is
\[
\hat H_{i,\mathrm{tot}}
=
\hat H_i+\hat V_i.
\]

Hence,
\begin{equation}
i\hbar
\frac{\partial}{\partial t}
\ket{\Phi_i(t)}
=
(\hat H_i+\hat V_i)
\ket{\Phi_i(t)}.
\tag{2}
\end{equation}

The free Hamiltonian describes the internal energy of the atom,
whereas \(\hat V_i\) represents the interaction between the atom
and the incident electromagnetic field.

\section{Free Hamiltonian of the Two-Level Atom}

The atomic energy states satisfy the energy eigenvalue equation
\[
\hat H_i\ket{\psi_n}
=
E_n\ket{\psi_n}.
\]

For the two states,
\[
\hat H_i\kete
=
E_e\kete,
\]
and
\[
\hat H_i\ketg
=
E_g\ketg.
\]

The energies of the two states are related to their angular
frequencies by
\[
E_e=\hbar\omega_e,
\qquad
E_g=\hbar\omega_g.
\]

Using the basis
\[
\kete=
\begin{bmatrix}
1\\
0
\end{bmatrix},
\qquad
\ketg=
\begin{bmatrix}
0\\
1
\end{bmatrix},
\]
write the free Hamiltonian in the general form
\[
\hat H_i
=
\begin{bmatrix}
a&b\\
c&d
\end{bmatrix}.
\]


% ============================================================
% PAGE 3
% ============================================================

Applying \(\hat H_i\) to the excited state gives
\[
\begin{aligned}
\hat H_i\kete
&=
\begin{bmatrix}
a&b\\
c&d
\end{bmatrix}
\begin{bmatrix}
1\\
0
\end{bmatrix}\\
&=
\begin{bmatrix}
a\\
c
\end{bmatrix}.
\end{aligned}
\]

From the energy eigenvalue equation,
\[
\hat H_i\kete
=
E_e\kete.
\]

Therefore,
\[
E_e\kete
=
E_e
\begin{bmatrix}
1\\
0
\end{bmatrix}
=
\begin{bmatrix}
E_e\\
0
\end{bmatrix}.
\]

Comparing the two expressions,
\[
\begin{bmatrix}
a\\
c
\end{bmatrix}
=
\begin{bmatrix}
E_e\\
0
\end{bmatrix},
\]
which gives
\[
a=E_e,
\qquad
c=0.
\]

Similarly, applying \(\hat H_i\) to the ground state,
\[
\begin{aligned}
\hat H_i\ketg
&=
\begin{bmatrix}
a&b\\
c&d
\end{bmatrix}
\begin{bmatrix}
0\\
1
\end{bmatrix}\\
&=
\begin{bmatrix}
b\\
d
\end{bmatrix}.
\end{aligned}
\]

The energy eigenvalue equation gives
\[
\hat H_i\ketg
=
E_g\ketg
=
E_g
\begin{bmatrix}
0\\
1
\end{bmatrix}
=
\begin{bmatrix}
0\\
E_g
\end{bmatrix}.
\]

Therefore,
\[
\begin{bmatrix}
b\\
d
\end{bmatrix}
=
\begin{bmatrix}
0\\
E_g
\end{bmatrix},
\]
and hence
\[
b=0,
\qquad
d=E_g.
\]

Thus,
\begin{equation}
\boxed{
\hat H_i
=
\begin{bmatrix}
E_e&0\\
0&E_g
\end{bmatrix}.
}
\tag{3}
\end{equation}

Using
\[
E_e=\hbar\omega_e,
\qquad
E_g=\hbar\omega_g,
\]
we obtain
\begin{equation}
\boxed{
\hat H_i
=
\begin{bmatrix}
\hbar\omega_e&0\\
0&\hbar\omega_g
\end{bmatrix}
=
\operatorname{diag}
(\hbar\omega_e,\hbar\omega_g).
}
\tag{4}
\end{equation}


% ============================================================
% PAGE 4
% ============================================================

\section{One-Dimensional Vapor-Cell Array Geometry}

Consider a one-dimensional array of vapor cells arranged along
the \(x\)-axis.

\begin{figure}[H]
    \centering
    \includegraphics[width=0.85\textwidth]{1D_array_geometry.png}
    \caption{1D uniform linear array geometry showing the inter-cell spacing $d$ and the position of the $i$-th vapor cell.}
    \label{fig:1D_array_geometry}
\end{figure}
Let the spacing between two adjacent vapor cells be \(d\), and
choose vapor cell \(1\) as the reference cell.

The position of vapor cell \(i\) is
\[
x_i=(i-1)d.
\]

Therefore,
\[
x_1=0,
\qquad
x_2=d,
\qquad
x_3=2d,
\qquad
\ldots
\]

The separation between the reference cell and cell \(i\) is
\[
x_i-x_1
=
(i-1)d.
\]

Consider user \(k\) and propagation path \(l\). Let
\(\theta_{l,k}\) denote the angle of arrival of this path measured
with respect to the array axis.

Under the far-field plane-wave assumption, the propagation paths
arriving at the different vapor cells are approximately parallel.

Hence,
\begin{equation}
\boxed{
\Delta r_{i,k,l}
=
(i-1)d\cos\theta_{l,k}.
}
\tag{5}
\end{equation}

where
\[
\Delta r_{i,k,l}
=
r_{i,k,l}-r_{1,k,l}.
\]

Here \(r_{i,k,l}\) denotes the propagation distance from user \(k\),
through path \(l\), to vapor cell \(i\).

For the second vapor cell,
\[
\Delta r_{2,k,l}
=
d\cos\theta_{l,k}.
\]

For the third vapor cell,
\[
\Delta r_{3,k,l}
=
2d\cos\theta_{l,k}.
\]

In general,
\[
\Delta r_{i,k,l}
=
(i-1)d\cos\theta_{l,k}.
\]


% ============================================================
% PAGE 5
% ============================================================

\section{Propagation Phase Difference}

Consider a sinusoidal plane wave represented by
\begin{equation}
E(r,t)
=
E_0
\cos
\left(
\omega t+kr+\phi_0
\right),
\tag{6}
\end{equation}
where
\[
k=\frac{2\pi}{\lambda}
\]
is the wavenumber.

If the propagation distance changes from \(r\) to
\(r+\Delta r\), the phase changes from
\[
kr
\]
to
\[
k(r+\Delta r).
\]

Therefore,
\[
\begin{aligned}
\Delta\phi
&=
k(r+\Delta r)-kr\\
&=
kr+k\Delta r-kr\\
&=
k\Delta r.
\end{aligned}
\]

Using
\[
k=\frac{2\pi}{\lambda},
\]
we obtain
\begin{equation}
\boxed{
\Delta\phi
=
\frac{2\pi}{\lambda}\Delta r.
}
\tag{7}
\end{equation}

For user \(k\), path \(l\), and vapor cell \(i\),
\[
\Delta r
=
\Delta r_{i,k,l}.
\]

Therefore,
\[
\begin{aligned}
\Delta\phi_{i,k,l}
&=
\frac{2\pi}{\lambda}
\Delta r_{i,k,l}\\
&=
\frac{2\pi}{\lambda}
(i-1)d\cos\theta_{l,k}\\
&=
(i-1)
\frac{2\pi d}{\lambda}
\cos\theta_{l,k}.
\end{aligned}
\]

Define the phase shift corresponding to one array spacing as
\begin{equation}
\boxed{
\phi_{l,k}
=
\frac{2\pi d}{\lambda}
\cos\theta_{l,k}.
}
\tag{8}
\end{equation}

Consequently,
\begin{equation}
\boxed{
\Delta\phi_{i,k,l}
=
(i-1)\phi_{l,k}.
}
\tag{9}
\end{equation}

For the reference vapor cell,
\[
i=1,
\]
and hence
\[
\Delta\phi_{1,k,l}=0.
\]

For the second vapor cell,
\[
i=2,
\]
and hence
\[
\Delta\phi_{2,k,l}=\phi_{l,k}.
\]

For the third vapor cell,
\[
i=3,
\]
and hence
\[
\Delta\phi_{3,k,l}=2\phi_{l,k}.
\]

Thus, the phase shift increases by \(\phi_{l,k}\) for every
additional array spacing.


% ============================================================
% PAGE 6
% ============================================================

\section{Electric Field of One Propagation Path}

Consider user \(k\) and propagation path \(l\).

The electric field corresponding to this propagation path can be
written as
\begin{equation}
\mathbf E_{k,l}(r,t)
=
\mathbf E_{0,k,l}
\cos
\left(
\omega t+kr+\phi_{0,k,l}
\right).
\tag{10}
\end{equation}

At vapor cell \(i\), the propagation distance is
\(r_{i,k,l}\). Therefore,
\begin{equation}
\mathbf E_{i,k,l}(t)
=
\mathbf E_{0,i,k,l}
\cos
\left(
\omega t
+
kr_{i,k,l}
+
\phi_{0,k,l}
\right).
\tag{11}
\end{equation}

Choose vapor cell \(1\) as the phase reference and define
\[
\Delta r_{i,k,l}
=
r_{i,k,l}-r_{1,k,l}.
\]

Therefore,
\[
r_{i,k,l}
=
r_{1,k,l}
+
\Delta r_{i,k,l}.
\]

Substituting this into (11),
\[
\begin{aligned}
\mathbf E_{i,k,l}(t)
&=
\mathbf E_{0,i,k,l}
\cos
\left[
\omega t
+
k
\left(
r_{1,k,l}+\Delta r_{i,k,l}
\right)
+
\phi_{0,k,l}
\right]\\
&=
\mathbf E_{0,i,k,l}
\cos
\left[
\omega t
+
k\Delta r_{i,k,l}
+
kr_{1,k,l}
+
\phi_{0,k,l}
\right].
\end{aligned}
\]

Define the common reference phase as
\[
\phi_{k,l}^{\mathrm{ref}}
=
kr_{1,k,l}
+
\phi_{0,k,l}.
\]

Hence,
\[
\mathbf E_{i,k,l}(t)
=
\mathbf E_{0,i,k,l}
\cos
\left[
\omega t
+
k\Delta r_{i,k,l}
+
\phi_{k,l}^{\mathrm{ref}}
\right].
\]

The phase \(\phi_{k,l}^{\mathrm{ref}}\) is common to all vapor cells, so it does not
change the relative phase across the array. It is absorbed into the complex
path coefficient \(\alpha_{l,k}\) introduced in Section~8. Without loss of
generality, therefore,
\[
\phi_{k,l}^{\mathrm{ref}}=0.
\]

Therefore,
\begin{equation}
\mathbf E_{i,k,l}(t)
=
\mathbf E_{0,i,k,l}
\cos
\left(
\omega t+k\Delta r_{i,k,l}
\right).
\tag{12}
\end{equation}

Using
\[
\Delta r_{i,k,l}
=
(i-1)d\cos\theta_{l,k},
\]
we obtain
\[
\begin{aligned}
\mathbf E_{i,k,l}(t)
&=
\mathbf E_{0,i,k,l}
\cos
\left[
\omega t
+
k(i-1)d\cos\theta_{l,k}
\right].
\end{aligned}
\]


% ============================================================
% PAGE 7
% ============================================================

Using
\[
k=\frac{2\pi}{\lambda},
\]
the electric field becomes
\[
\begin{aligned}
\mathbf E_{i,k,l}(t)
&=
\mathbf E_{0,i,k,l}
\cos
\left[
\omega t
+
(i-1)
\frac{2\pi d}{\lambda}
\cos\theta_{l,k}
\right].
\end{aligned}
\]

From (8),
\[
\phi_{l,k}
=
\frac{2\pi d}{\lambda}
\cos\theta_{l,k}.
\]

Therefore,
\begin{equation}
\boxed{
\mathbf E_{i,k,l}(t)
=
\mathbf E_{0,i,k,l}
\cos
\left[
\omega t+(i-1)\phi_{l,k}
\right].
}
\tag{13}
\end{equation}

\section{Vector Representation and Polarization}

The electric field is a vector quantity. Its amplitude can therefore
be separated into a scalar amplitude and a polarization direction.

Let
\[
\boldsymbol{\epsilon}_{i,k,l}
\]
denote the polarization vector corresponding to user \(k\),
path \(l\), and vapor cell \(i\).

The vector amplitude can be written as
\[
\mathbf E_{0,i,k,l}
=
E_{0,i,k,l}
\boldsymbol{\epsilon}_{i,k,l}.
\]

Substituting this into (13),
\[
\begin{aligned}
\mathbf E_{i,k,l}(t)
&=
E_{0,i,k,l}
\boldsymbol{\epsilon}_{i,k,l}
\cos
\left[
\omega t+(i-1)\phi_{l,k}
\right].
\end{aligned}
\]

Hence,
\begin{equation}
\boxed{
\mathbf E_{i,k,l}(t)
=
E_{0,i,k,l}
\boldsymbol{\epsilon}_{i,k,l}
\cos
\left[
\omega t+(i-1)\phi_{l,k}
\right].
}
\tag{14}
\end{equation}

The polarization vector represents the direction of the electric
field, while \(E_{0,i,k,l}\) represents its scalar amplitude.


% ============================================================
% PAGE 8
% ============================================================

\section{Transmitted Symbol and Propagation Coefficient}

Consider user \(k\) transmitting the symbol
\[
s_k.
\]

Let
\[
\alpha_{l,k}
\]
denote the complex coefficient associated with propagation path
\(l\) from user \(k\).

The scalar field amplitude is represented as
\[
E_{0,i,k,l}
=
\alpha_{l,k}s_k.
\]

Substituting this relation into (14),
\[
\begin{aligned}
\mathbf E_{i,k,l}(t)
&=
\alpha_{l,k}s_k
\boldsymbol{\epsilon}_{i,k,l}
\cos
\left[
\omega t+(i-1)\phi_{l,k}
\right].
\end{aligned}
\]

Therefore,
\begin{equation}
\boxed{
\mathbf E_{i,k,l}(t)
=
\boldsymbol{\epsilon}_{i,k,l}
\alpha_{l,k}s_k
\cos
\left[
\omega t+(i-1)\phi_{l,k}
\right].
}
\tag{15}
\end{equation}

The three terms in this expression represent the polarization
direction, propagation-path coefficient, and transmitted symbol,
respectively:
\[
\boldsymbol{\epsilon}_{i,k,l},
\qquad
\alpha_{l,k},
\qquad
s_k.
\]

The phase accumulated at vapor cell \(i\) is represented by
\[
(i-1)\phi_{l,k}.
\]


% ============================================================
% PAGE 9
% ============================================================

\section{Superposition of All Propagation Paths}

For user \(k\), the transmitted signal reaches the vapor cell
through \(L_k\) propagation paths.

The total electric field produced by user \(k\) is therefore the
sum of the fields associated with all its propagation paths:
\[
\mathbf E_{i,k}(t)
=
\sum_{l=1}^{L_k}
\mathbf E_{i,k,l}(t).
\]

Substituting (15),
\[
\begin{aligned}
\mathbf E_{i,k}(t)
&=
\sum_{l=1}^{L_k}
\boldsymbol{\epsilon}_{i,k,l}
\alpha_{l,k}s_k
\cos
\left[
\omega t+(i-1)\phi_{l,k}
\right].
\end{aligned}
\]

The vapor cell receives signals from all \(K\) users. Therefore,
the total electric field at vapor cell \(i\) is
\[
\mathbf E_i(t)
=
\sum_{k=1}^{K}
\mathbf E_{i,k}(t).
\]

Substituting the expression above,
\[
\begin{aligned}
\mathbf E_i(t)
&=
\sum_{k=1}^{K}
\sum_{l=1}^{L_k}
\boldsymbol{\epsilon}_{i,k,l}
\alpha_{l,k}s_k
\cos
\left[
\omega t+(i-1)\phi_{l,k}
\right].
\end{aligned}
\]

Thus,
\begin{equation}
\boxed{
\mathbf E_i(t)
=
\sum_{k=1}^{K}
\sum_{l=1}^{L_k}
\boldsymbol{\epsilon}_{i,k,l}
\alpha_{l,k}s_k
\cos
\left[
\omega t+(i-1)\phi_{l,k}
\right].
}
\tag{16}
\end{equation}


% ============================================================
% PAGE 10
% ============================================================

\section{Electric-Dipole Interaction}

The electromagnetic field interacts with the electric dipole of
the atom.

For an electric dipole moment \(\boldsymbol{\mu}\) in an electric
field \(\mathbf E\), the interaction energy is represented by the
dipole--field product.

For a charge \(q\) located at position \(\mathbf r\), the electric
dipole moment is
\[
\boldsymbol{\mu}
=
q\mathbf r.
\]

Therefore, the interaction energy can be written in terms of the
position vector and electric field as
\[
V
=
q\mathbf r^{T}\mathbf E.
\]

In quantum mechanics, the position of the electron is represented
by a position operator. Hence,
\[
\hat{\mathbf r}
=
\begin{bmatrix}
\hat x\\
\hat y\\
\hat z
\end{bmatrix}.
\]

The interaction Hamiltonian for the atom in vapor cell \(i\) is
therefore written as
\begin{equation}
\boxed{
\hat V_i
=
q\hat{\mathbf r}^{\,T}\mathbf E_i(t).
}
\tag{17}
\end{equation}

Physically, the dipole interaction energy is \(-\boldsymbol{\mu}^{T}\mathbf E\).
The sign convention in (17) follows Cui \emph{et al.}; it only reverses the sign
of every channel and reference coefficient simultaneously, so all magnitudes,
and hence all later results, are unchanged.

The electric field is a vector,
\[
\mathbf E_i(t)
=
\begin{bmatrix}
E_{i,x}(t)\\
E_{i,y}(t)\\
E_{i,z}(t)
\end{bmatrix},
\]
and hence
\[
\begin{aligned}
\hat{\mathbf r}^{\,T}\mathbf E_i(t)
&=
\begin{bmatrix}
\hat x&\hat y&\hat z
\end{bmatrix}
\begin{bmatrix}
E_{i,x}(t)\\
E_{i,y}(t)\\
E_{i,z}(t)
\end{bmatrix}\\
&=
\hat xE_{i,x}(t)
+
\hat yE_{i,y}(t)
+
\hat zE_{i,z}(t).
\end{aligned}
\]

Therefore,
\[
\boxed{
\hat V_i
=
q
\left[
\hat xE_{i,x}(t)
+
\hat yE_{i,y}(t)
+
\hat zE_{i,z}(t)
\right].
}
\tag{18}
\]

At this stage, the interaction Hamiltonian is expressed in terms
of the position operator. The next step is to represent this
operator in the two-level basis
\[
\{\ketg,\kete\}.
\]



% ============================================================
% END OF FIRST 10 SCANNED PAGES
%
% Last equation number: (18)
% Next block should continue from equation (19).
% ===========================================================

% ============================================================
% PAGES 11--20 OF THE SCANNED DERIVATION
% CONTINUATION FROM EQUATION (18)
% ============================================================

\section{Two-Level Representation of the Position Operator}

In quantum mechanics, the position of the electron is represented by
a position operator rather than an ordinary definite coordinate.

The position operator is written as
\[
\hat{\mathbf r}
=
\begin{bmatrix}
\hat x\\
\hat y\\
\hat z
\end{bmatrix}.
\]

For the two-level Rydberg-atom system, the state can be expressed as
\[
\ket{\Phi_i(t)}
=
\zeta_{g,i}(t)\ketg
+
\zeta_{e,i}(t)\kete.
\]

The completeness relation for the two-level basis is
\begin{equation}
\boxed{
\hat I
=
\ketg\bra g
+
\kete\bra e.
}
\tag{19}
\end{equation}

Using the completeness relation on both sides of the position
operator,
\[
\hat{\mathbf r}
=
\hat I\hat{\mathbf r}\hat I.
\]

Therefore,
\[
\begin{aligned}
\hat{\mathbf r}
&=
\left(
\ketg\bra g+\kete\bra e
\right)
\hat{\mathbf r}
\left(
\ketg\bra g+\kete\bra e
\right).
\end{aligned}
\]

Expanding the product gives
\[
\begin{aligned}
\hat{\mathbf r}
={}&
\ketg\bra g
\hat{\mathbf r}
\ketg\bra g
+
\ketg\bra g
\hat{\mathbf r}
\kete\bra e\\
&+
\kete\bra e
\hat{\mathbf r}
\ketg\bra g
+
\kete\bra e
\hat{\mathbf r}
\kete\bra e.
\end{aligned}
\]

The diagonal matrix elements are
\[
\bra g\hat{\mathbf r}\ketg
\]
and
\[
\bra e\hat{\mathbf r}\kete,
\]
which represent the position-operator expectation values within
the ground and excited states.

The off-diagonal matrix elements
\[
\bra e\hat{\mathbf r}\ketg
\]
and
\[
\bra g\hat{\mathbf r}\kete
\]
are the transition matrix elements between the two atomic states.

The position operator is odd under parity, and the two dipole-coupled states
\(\ketg\) and \(\kete\) have opposite parity. Therefore the diagonal elements
vanish,
\[
\bra g\hat{\mathbf r}\ketg=\bra e\hat{\mathbf r}\kete=\mathbf 0,
\]
and only the transition terms remain. Hence,
\begin{equation}
\boxed{
\hat{\mathbf r}
=
\bra e\hat{\mathbf r}\ketg
\kete\bra g
+
\bra g\hat{\mathbf r}\kete
\ketg\bra e.
}
\tag{20}
\end{equation}


\section{Transition Dipole Moment}

The interaction Hamiltonian obtained previously is
\[
\hat V_i
=
q\hat{\mathbf r}^{\,T}\mathbf E_i(t).
\]

Substituting the two-level representation of the position operator,
\[
\begin{aligned}
\hat V_i
&=
q
\left[
\bra e\hat{\mathbf r}\ketg
\kete\bra g
+
\bra g\hat{\mathbf r}\kete
\ketg\bra e
\right]^{T}
\mathbf E_i(t).
\end{aligned}
\]

The electric field is an ordinary classical vector. Therefore, it
can be included in the corresponding scalar products.

Define the transition dipole moment between the ground and excited
states as
\begin{equation}
\boxed{
\boldsymbol{\mu}_{eg}
=
q\bra e\hat{\mathbf r}\ketg.
}
\tag{21}
\end{equation}

Similarly, define
\begin{equation}
\boxed{
\boldsymbol{\mu}_{ge}
=
q\bra g\hat{\mathbf r}\kete.
}
\tag{22}
\end{equation}

Therefore,
\[
q
\bra e\hat{\mathbf r}\ketg
=
\boldsymbol{\mu}_{eg},
\]
and
\[
q
\bra g\hat{\mathbf r}\kete
=
\boldsymbol{\mu}_{ge}.
\]

Substituting these definitions into the interaction Hamiltonian,
\[
\hat V_i
=
\left(
\boldsymbol{\mu}_{eg}^{T}
\mathbf E_i(t)
\right)
\kete\bra g
+
\left(
\boldsymbol{\mu}_{ge}^{T}
\mathbf E_i(t)
\right)
\ketg\bra e.
\]

The position operator is Hermitian:
\[
\hat{\mathbf r}^{\dagger}
=
\hat{\mathbf r}.
\]

Hence,
\[
\left(
\bra e\hat{\mathbf r}\ketg
\right)^{*}
=
\bra g\hat{\mathbf r}\kete.
\]

Therefore,
\[
\boldsymbol{\mu}_{ge}
=
\boldsymbol{\mu}_{eg}^{*}.
\]

Consequently,
\[
\left(
\boldsymbol{\mu}_{eg}^{T}\mathbf E_i(t)
\right)^{*}
=
\boldsymbol{\mu}_{ge}^{T}\mathbf E_i(t),
\]
since the physical electric field is real. Writing the complex amplitude
\(\alpha_{l,k}s_k\) inside the cosine in (15) is shorthand for the real part of
the corresponding phasor; because \(\boldsymbol\mu_{eg}\) and the polarization
vectors can be taken real, this does not change \(|C_i|\) or any later result.

Define
\begin{equation}
\boxed{
\bar{\omega}_i(t)
=
\boldsymbol{\mu}_{eg}^{T}\mathbf E_i(t).
}
\tag{23}
\end{equation}

Then,
\[
\bar{\omega}_i^{*}(t)
=
\boldsymbol{\mu}_{ge}^{T}\mathbf E_i(t).
\]

Therefore, the interaction Hamiltonian becomes
\begin{equation}
\boxed{
\hat V_i
=
\bar{\omega}_i(t)
\kete\bra g
+
\bar{\omega}_i^{*}(t)
\ketg\bra e.
}
\tag{24}
\end{equation}


\section{Matrix Representation of the Interaction Hamiltonian}

Using
\[
\kete
=
\begin{bmatrix}
1\\
0
\end{bmatrix},
\qquad
\ketg
=
\begin{bmatrix}
0\\
1
\end{bmatrix},
\]
we obtain
\[
\kete\bra g
=
\begin{bmatrix}
1\\
0
\end{bmatrix}
\begin{bmatrix}
0&1
\end{bmatrix}
=
\begin{bmatrix}
0&1\\
0&0
\end{bmatrix}.
\]

Similarly,
\[
\ketg\bra e
=
\begin{bmatrix}
0\\
1
\end{bmatrix}
\begin{bmatrix}
1&0
\end{bmatrix}
=
\begin{bmatrix}
0&0\\
1&0
\end{bmatrix}.
\]

Substituting these into (24),
\[
\begin{aligned}
\hat V_i
&=
\bar{\omega}_i(t)
\begin{bmatrix}
0&1\\
0&0
\end{bmatrix}
+
\bar{\omega}_i^{*}(t)
\begin{bmatrix}
0&0\\
1&0
\end{bmatrix}\\
&=
\begin{bmatrix}
0&\bar{\omega}_i(t)\\
\bar{\omega}_i^{*}(t)&0
\end{bmatrix}.
\end{aligned}
\]

Hence,
\begin{equation}
\boxed{
\hat V_i
=
\begin{bmatrix}
0&\bar{\omega}_i(t)\\
\bar{\omega}_i^{*}(t)&0
\end{bmatrix}.
}
\tag{25}
\end{equation}


\section{Interaction Coefficient in Terms of the Electric Field}

From the definition
\[
\bar{\omega}_i(t)
=
\boldsymbol{\mu}_{eg}^{T}\mathbf E_i(t),
\]
and the total electric field obtained previously,
\[
\mathbf E_i(t)
=
\sum_{k=1}^{K}
\sum_{l=1}^{L_k}
\boldsymbol{\epsilon}_{i,k,l}
\alpha_{l,k}s_k
\cos
\left[
\omega t+(i-1)\phi_{l,k}
\right],
\]
we obtain
\[
\begin{aligned}
\bar{\omega}_i(t)
&=
\boldsymbol{\mu}_{eg}^{T}
\sum_{k=1}^{K}
\sum_{l=1}^{L_k}
\boldsymbol{\epsilon}_{i,k,l}
\alpha_{l,k}s_k
\cos
\left[
\omega t+(i-1)\phi_{l,k}
\right].
\end{aligned}
\]

Thus,
\begin{equation}
\boxed{
\bar{\omega}_i(t)
=
\sum_{k=1}^{K}
\sum_{l=1}^{L_k}
\boldsymbol{\mu}_{eg}^{T}
\boldsymbol{\epsilon}_{i,k,l}
\alpha_{l,k}s_k
\cos
\left[
\omega t+(i-1)\phi_{l,k}
\right].
}
\tag{26}
\end{equation}

Similarly,
\[
\begin{aligned}
\bar{\omega}_i^{*}(t)
=
\sum_{k=1}^{K}
\sum_{l=1}^{L_k}
\boldsymbol{\mu}_{ge}^{T}
\boldsymbol{\epsilon}_{i,k,l}
\alpha_{l,k}^{*}s_k^{*}
\cos
\left[
\omega t+(i-1)\phi_{l,k}
\right].
\end{aligned}
\]


\section{Total Hamiltonian}

The free Hamiltonian is
\[
\hat H_i
=
\begin{bmatrix}
\hbar\omega_e&0\\
0&\hbar\omega_g
\end{bmatrix},
\]
and the interaction Hamiltonian is
\[
\hat V_i
=
\begin{bmatrix}
0&\bar{\omega}_i(t)\\
\bar{\omega}_i^{*}(t)&0
\end{bmatrix}.
\]

Therefore,
\[
\begin{aligned}
\hat H_i+\hat V_i
&=
\begin{bmatrix}
\hbar\omega_e&0\\
0&\hbar\omega_g
\end{bmatrix}
+
\begin{bmatrix}
0&\bar{\omega}_i(t)\\
\bar{\omega}_i^{*}(t)&0
\end{bmatrix}\\
&=
\begin{bmatrix}
\hbar\omega_e&\bar{\omega}_i(t)\\
\bar{\omega}_i^{*}(t)&\hbar\omega_g
\end{bmatrix}.
\end{aligned}
\]

Hence,
\begin{equation}
\boxed{
\hat H_i+\hat V_i
=
\begin{bmatrix}
\hbar\omega_e&\bar{\omega}_i(t)\\
\bar{\omega}_i^{*}(t)&\hbar\omega_g
\end{bmatrix}.
}
\tag{27}
\end{equation}


\section{Coupled Schr\"odinger Equations}

Using
\[
\ket{\Phi_i(t)}
=
\begin{bmatrix}
\zeta_{e,i}(t)\\
\zeta_{g,i}(t)
\end{bmatrix},
\]
the Schr\"odinger equation becomes
\[
i\hbar
\frac{d}{dt}
\begin{bmatrix}
\zeta_{e,i}(t)\\
\zeta_{g,i}(t)
\end{bmatrix}
=
\begin{bmatrix}
\hbar\omega_e&\bar{\omega}_i(t)\\
\bar{\omega}_i^{*}(t)&\hbar\omega_g
\end{bmatrix}
\begin{bmatrix}
\zeta_{e,i}(t)\\
\zeta_{g,i}(t)
\end{bmatrix}.
\]

Multiplying the matrices,
\[
i\hbar
\begin{bmatrix}
\dot{\zeta}_{e,i}(t)\\
\dot{\zeta}_{g,i}(t)
\end{bmatrix}
=
\begin{bmatrix}
\hbar\omega_e\zeta_{e,i}(t)
+
\bar{\omega}_i(t)\zeta_{g,i}(t)
\\[4pt]
\bar{\omega}_i^{*}(t)\zeta_{e,i}(t)
+
\hbar\omega_g\zeta_{g,i}(t)
\end{bmatrix}.
\]

Therefore,
\begin{equation}
\boxed{
i\hbar\dot{\zeta}_{e,i}
=
\hbar\omega_e\zeta_{e,i}
+
\bar{\omega}_i(t)\zeta_{g,i}.
}
\tag{28a}
\end{equation}

and
\begin{equation}
\boxed{
i\hbar\dot{\zeta}_{g,i}
=
\bar{\omega}_i^{*}(t)\zeta_{e,i}
+
\hbar\omega_g\zeta_{g,i}.
}
\tag{28b}
\end{equation}

The terms
\[
\hbar\omega_e\zeta_{e,i}
\]
and
\[
\hbar\omega_g\zeta_{g,i}
\]
represent the free atomic evolution, while the terms containing
\(\bar{\omega}_i(t)\) represent the interaction with the
electromagnetic field.


\section{Removal of the Free-Hamiltonian Evolution}

We start from the time-dependent Schr\"odinger equation for the
free atomic evolution,
\[
i\hbar
\frac{\partial}{\partial t}
|\Phi_i(t)\rangle
=
\hat H_i|\Phi_i(t)\rangle.
\]

For the free Hamiltonian,
\[
\hat H_i
=
\begin{bmatrix}
\hbar\omega_e & 0\\
0 & \hbar\omega_g
\end{bmatrix}.
\]

Consider first the excited-state component. From the free
Schr\"odinger equation,
\[
i\hbar
\frac{d\zeta_{e,i}(t)}{dt}
=
\hbar\omega_e\zeta_{e,i}(t).
\]

we obtain
\[
\frac{d\zeta_{e,i}(t)}{dt}
=
-j\omega_e\zeta_{e,i}(t).
\]

Now separate the variables:
\[
\frac{1}{\zeta_{e,i}(t)}
\,d\zeta_{e,i}(t)
=
-j\omega_e\,dt.
\]

Integrating both sides,
\[
\int
\frac{1}{\zeta_{e,i}}
\,d\zeta_{e,i}
=
\int
-j\omega_e\,dt,
\]
which gives
\[
\ln\zeta_{e,i}(t)
=
-j\omega_e t+C.
\]

Exponentiating,
\[
\zeta_{e,i}(t)
=
Ce^{-j\omega_e t}.
\]

Thus, the free evolution of the excited-state component contains
the phase factor
\[
\boxed{
e^{-j\omega_e t}
}.
\]

Similarly, for the ground-state component,
\[
i\hbar
\frac{d\zeta_{g,i}(t)}{dt}
=
\hbar\omega_g\zeta_{g,i}(t),
\]
which gives
\[
\frac{d\zeta_{g,i}(t)}{dt}
=
-j\omega_g\zeta_{g,i}(t),
\]
and hence
\[
\boxed{
e^{-j\omega_g t}
}
\]
for the ground state.

We now define the slowly varying amplitudes \(a_i(t)\) and
\(b_i(t)\) by
\begin{equation}
\boxed{
\zeta_{e,i}(t)
=
a_i(t)e^{-j\omega_e t}.
}
\tag{29a}
\end{equation}

Similarly,
\begin{equation}
\boxed{
\zeta_{g,i}(t)
=
b_i(t)e^{-j\omega_g t}.
}
\tag{29b}
\end{equation}

The functions \(a_i(t)\) and \(b_i(t)\) contain the evolution
produced by the electromagnetic interaction after the free atomic
oscillation has been removed.

For the excited state,
\[
\zeta_{e,i}(t)
=
a_i(t)e^{-j\omega_e t}.
\]

Differentiating,
\[
\begin{aligned}
\dot{\zeta}_{e,i}(t)
&=
\frac{d}{dt}
\left[
a_i(t)e^{-j\omega_e t}
\right]\\
&=
\dot a_i(t)e^{-j\omega_e t}
+
a_i(t)
\frac{d}{dt}
\left(
e^{-j\omega_e t}
\right)\\
&=
\dot a_i(t)e^{-j\omega_e t}
-
j\omega_ea_i(t)e^{-j\omega_e t}\\
&=
\left(
\dot a_i-j\omega_ea_i
\right)
e^{-j\omega_e t}.
\end{aligned}
\]

Hence,
\begin{equation}
\boxed{
\dot{\zeta}_{e,i}
=
\left(
\dot a_i-j\omega_ea_i
\right)
e^{-j\omega_e t}.
}
\tag{30}
\end{equation}

Similarly,
\[
\begin{aligned}
\dot{\zeta}_{g,i}
&=
\frac{d}{dt}
\left[
b_i(t)e^{-j\omega_g t}
\right]\\
&=
\dot b_i(t)e^{-j\omega_g t}
-
j\omega_gb_i(t)e^{-j\omega_g t}\\
&=
\left(
\dot b_i-j\omega_gb_i
\right)
e^{-j\omega_g t}.
\end{aligned}
\]

Therefore,
\begin{equation}
\boxed{
\dot{\zeta}_{g,i}
=
\left(
\dot b_i-j\omega_gb_i
\right)
e^{-j\omega_g t}.
}
\tag{31}
\end{equation}

\section{Transformation of the Excited-State Equation}

Substitute
\[
\zeta_{e,i}
=
a_i e^{-j\omega_e t},
\qquad
\zeta_{g,i}
=
b_i e^{-j\omega_g t},
\]
and (30) into (28a):
\[
i\hbar
\left(
\dot a_i-j\omega_ea_i
\right)
e^{-j\omega_e t}
=
\hbar\omega_e
a_i e^{-j\omega_e t}
+
\bar{\omega}_i(t)
b_i e^{-j\omega_g t}.
\]

Expanding the left-hand side,
\[
i\hbar\dot a_i e^{-j\omega_e t}
+
\hbar\omega_ea_i e^{-j\omega_e t}
=
\hbar\omega_ea_i e^{-j\omega_e t}
+
\bar{\omega}_i(t)b_i e^{-j\omega_g t}.
\]

The free-Hamiltonian terms cancel:
\[
i\hbar\dot a_i e^{-j\omega_e t}
=
\bar{\omega}_i(t)b_i e^{-j\omega_g t}.
\]

Multiplying both sides by
\[
e^{j\omega_e t},
\]
we obtain
\[
i\hbar\dot a_i
=
\bar{\omega}_i(t)b_i
e^{j(\omega_e-\omega_g)t}.
\]

Define the atomic transition angular frequency as
\begin{equation}
\boxed{
\omega_{eg}
=
\omega_e-\omega_g.
}
\tag{32}
\end{equation}

Hence,
\begin{equation}
\boxed{
i\hbar\dot a_i
=
\bar{\omega}_i(t)
b_i
e^{j\omega_{eg}t}.
}
\tag{33}
\end{equation}

Similarly, substituting the transformations into the ground-state
equation gives
\begin{equation}
\boxed{
i\hbar\dot b_i
=
\bar{\omega}_i^{*}(t)
a_i
e^{-j\omega_{eg}t}.
}
\tag{34}
\end{equation}


\section{Substitution of the Electromagnetic Interaction}

From (26),
\[
\bar{\omega}_i(t)
=
\sum_{k=1}^{K}
\sum_{l=1}^{L_k}
\boldsymbol{\mu}_{eg}^{T}
\boldsymbol{\epsilon}_{i,k,l}
\alpha_{l,k}s_k
\cos
\left[
\omega t+(i-1)\phi_{l,k}
\right].
\]

Define
\begin{equation}
\boxed{
A_{i,k,l}
=
\boldsymbol{\mu}_{eg}^{T}
\boldsymbol{\epsilon}_{i,k,l}
\alpha_{l,k}s_k.
}
\tag{35}
\end{equation}

Then,
\[
\bar{\omega}_i(t)
=
\sum_{k=1}^{K}
\sum_{l=1}^{L_k}
A_{i,k,l}
\cos
\left[
\omega t+(i-1)\phi_{l,k}
\right].
\]

Substituting this expression into (33),
\[
\begin{aligned}
i\hbar\dot a_i
=
\sum_{k=1}^{K}
\sum_{l=1}^{L_k}
&
A_{i,k,l}
\cos
\left[
\omega t+(i-1)\phi_{l,k}
\right]\\
&\times
e^{j\omega_{eg}t}
b_i.
\end{aligned}
\]

Using the identity
\[
\cos x
=
\frac{e^{jx}+e^{-jx}}{2},
\]
we obtain
\[
\begin{aligned}
\cos
\left[
\omega t+(i-1)\phi_{l,k}
\right]
=
\frac{1}{2}
\Big[
&
e^{j[\omega t+(i-1)\phi_{l,k}]}\\
&+
e^{-j[\omega t+(i-1)\phi_{l,k}]}
\Big].
\end{aligned}
\]

Therefore,
\[
\begin{aligned}
i\hbar\dot a_i
=
\frac{1}{2}
\sum_{k=1}^{K}
\sum_{l=1}^{L_k}
A_{i,k,l}
\Big[
&
e^{j[\omega t+(i-1)\phi_{l,k}]}\\
&+
e^{-j[\omega t+(i-1)\phi_{l,k}]}
\Big]
e^{j\omega_{eg}t}
b_i.
\end{aligned}
\]

Combining the exponential terms,
\[
\begin{aligned}
i\hbar\dot a_i
=
\frac{1}{2}
\sum_{k=1}^{K}
\sum_{l=1}^{L_k}
A_{i,k,l}
\Big[
&
e^{j(\omega+\omega_{eg})t}
e^{j(i-1)\phi_{l,k}}
\\
&+
e^{j(\omega_{eg}-\omega)t}
e^{-j(i-1)\phi_{l,k}}
\Big]
b_i.
\end{aligned}
\tag{36}
\]


\section{Resonance Condition}

The electromagnetic field is considered resonant with the atomic
transition when
\[
\boxed{
\omega=\omega_{eg}.
}
\tag{37}
\]

Therefore,
\[
\omega+\omega_{eg}=2\omega,
\]
and
\[
\omega_{eg}-\omega=0.
\]

Substituting these relations into (36),
\[
\begin{aligned}
i\hbar\dot a_i
=
\frac{1}{2}
\sum_{k=1}^{K}
\sum_{l=1}^{L_k}
A_{i,k,l}
\Big[
&
e^{j2\omega t}
e^{j(i-1)\phi_{l,k}}
\\
&+
e^{-j(i-1)\phi_{l,k}}
\Big]
b_i.
\end{aligned}
\tag{38}
\]


\section{Rotating-Wave Approximation}

Under the rotating-wave approximation, terms oscillating rapidly
compared with the atomic transition dynamics are neglected.

The first term in (38),
\[
e^{j2\omega t}
e^{j(i-1)\phi_{l,k}},
\]
contains the rapidly oscillating factor
\[
e^{j2\omega t}.
\]

The second term,
\[
e^{-j(i-1)\phi_{l,k}},
\]
is slowly varying with respect to the atomic transition dynamics.

Therefore, the rapidly oscillating term is neglected, giving
\[
\boxed{
i\hbar\dot a_i
=
\frac{1}{2}
\sum_{k=1}^{K}
\sum_{l=1}^{L_k}
A_{i,k,l}
e^{-j(i-1)\phi_{l,k}}
b_i.
}
\tag{39}
\]

Substituting the definition of \(A_{i,k,l}\),
\[
\begin{aligned}
i\hbar\dot a_i
=
\frac{1}{2}
\sum_{k=1}^{K}
\sum_{l=1}^{L_k}
&
\boldsymbol{\mu}_{eg}^{T}
\boldsymbol{\epsilon}_{i,k,l}
\alpha_{l,k}s_k\\
&\times
e^{-j(i-1)\phi_{l,k}}
b_i.
\end{aligned}
\]

Define
\begin{equation}
\boxed{
C_i
=
\sum_{k=1}^{K}
\sum_{l=1}^{L_k}
\boldsymbol{\mu}_{eg}^{T}
\boldsymbol{\epsilon}_{i,k,l}
\alpha_{l,k}s_k
e^{-j(i-1)\phi_{l,k}}.
}
\tag{40}
\end{equation}

Then,
\begin{equation}
\boxed{
i\hbar\dot a_i
=
\frac{C_i}{2}b_i.
}
\tag{41}
\end{equation}

Similarly, taking the conjugate interaction coefficient in the
ground-state equation,
\begin{equation}
\boxed{
i\hbar\dot b_i
=
\frac{C_i^{*}}{2}a_i.
}
\tag{42}
\end{equation}


\section{Second-Order Differential Equation}

From (41),
\[
i\hbar\dot a_i
=
\frac{C_i}{2}b_i.
\]

Since \(C_i\) is constant with respect to time in the resonant
RWA model,
\[
i\hbar\ddot a_i
=
\frac{C_i}{2}\dot b_i.
\]

From (42),
\[
i\hbar\dot b_i
=
\frac{C_i^{*}}{2}a_i.
\]

Therefore,
\[
\dot b_i
=
\frac{C_i^{*}}{2i\hbar}a_i.
\]

Using
\[
\frac{1}{i}=-j,
\]
we obtain
\[
\dot b_i
=
-\frac{jC_i^{*}}{2\hbar}a_i.
\]

Substituting this into the differentiated form of (41),
\[
\begin{aligned}
i\hbar\ddot a_i
&=
\frac{C_i}{2}
\left(
-\frac{jC_i^{*}}{2\hbar}a_i
\right)\\
&=
-\frac{jC_iC_i^{*}}{4\hbar}a_i.
\end{aligned}
\]

Since
\[
C_iC_i^{*}
=
|C_i|^2,
\]
we have
\[
i\hbar\ddot a_i
=
-\frac{j|C_i|^2}{4\hbar}a_i.
\]

Dividing by \(i\hbar\),
\[
\ddot a_i
=
-\frac{|C_i|^2}{4\hbar^2}a_i.
\]

Therefore,
\begin{equation}
\boxed{
\ddot a_i
+
\frac{|C_i|^2}{4\hbar^2}a_i
=
0.
}
\tag{43}
\end{equation}

\section{Solution of the Excited-State Dynamics}

The second-order differential equation obtained in (43) is

\begin{equation}
\ddot{a}_i(t)
+
\frac{|C_i|^2}{4\hbar^2}a_i(t)
=0.
\tag{44}
\end{equation}

Its general solution is

\begin{equation}
a_i(t)
=
A\cos\left(\frac{|C_i|t}{2\hbar}\right)
+
B\sin\left(\frac{|C_i|t}{2\hbar}\right).
\tag{45}
\end{equation}

The atom is initially prepared in the ground state. Therefore,

\begin{equation}
\zeta_{e,i}(0)=0,
\qquad
\zeta_{g,i}(0)=1.
\tag{46}
\end{equation}

Since

\begin{equation}
\zeta_{e,i}(t)=a_i(t)e^{-j\omega_e t},
\tag{47}
\end{equation}

and $e^{-j\omega_e 0}=1$, the initial excited-state condition gives

\begin{equation}
a_i(0)=0.
\tag{48}
\end{equation}

Evaluating (45) at $t=0$ gives

\begin{equation}
a_i(0)=A.
\tag{49}
\end{equation}

Hence,

\begin{equation}
A=0.
\tag{50}
\end{equation}

Therefore, the transformed excited-state amplitude becomes

\begin{equation}
a_i(t)
=
B\sin\left(\frac{|C_i|t}{2\hbar}\right).
\tag{51}
\end{equation}

The coefficient $B$ is determined from the first-order equation obtained from the rotating-wave approximation,

\begin{equation}
i\hbar\dot a_i(t)=\frac{C_i}{2}b_i(t).
\tag{52}
\end{equation}

At $t=0$, $b_i(0)=1$, and therefore

\begin{equation}
i\hbar\dot a_i(0)=\frac{C_i}{2}.
\tag{53}
\end{equation}

Differentiating (51) gives

\begin{equation}
\dot a_i(t)
=
B\frac{|C_i|}{2\hbar}
\cos\left(\frac{|C_i|t}{2\hbar}\right).
\tag{54}
\end{equation}

Thus, at $t=0$,

\begin{equation}
\dot a_i(0)=B\frac{|C_i|}{2\hbar}.
\tag{55}
\end{equation}

Substituting (55) into (53) gives

\begin{equation}
iB|C_i|=C_i.
\tag{56}
\end{equation}

For $C_i\neq0$,

\begin{equation}
B=-j\frac{C_i}{|C_i|}.
\tag{57}
\end{equation}

Hence,

\begin{equation}
a_i(t)
=
-j\frac{C_i}{|C_i|}
\sin\left(\frac{|C_i|t}{2\hbar}\right).
\tag{58}
\end{equation}

Using (47), the original excited-state amplitude is

\begin{equation}
\zeta_{e,i}(t)
=
-j\frac{C_i}{|C_i|}
\sin\left(\frac{|C_i|t}{2\hbar}\right)e^{-j\omega_e t}.
\tag{59}
\end{equation}

The excited-state probability is defined as the squared magnitude of the excited-state probability amplitude,

\begin{equation}
P_e(t)=|\zeta_{e,i}(t)|^2.
\tag{60}
\end{equation}

Using the multiplicative property of the complex magnitude,

\begin{equation}
|abc|^2=|a|^2|b|^2|c|^2,
\tag{61}
\end{equation}

and the identities

\begin{equation}
|-j|^2=1,
\qquad
\left|\frac{C_i}{|C_i|}\right|^2=1,
\qquad
|e^{-j\omega_e t}|^2=1,
\tag{62}
\end{equation}

we obtain

\begin{equation}
P_e(t)
=
\sin^2\left(\frac{|C_i|t}{2\hbar}\right).
\tag{63}
\end{equation}

Define the Rabi frequency by

\begin{equation}
\boxed{
\Omega_i=\frac{|C_i|}{\hbar}
}.
\tag{64}
\end{equation}

Therefore,

\begin{equation}
\frac{|C_i|t}{2\hbar}=\frac{\Omega_i t}{2}.
\tag{65}
\end{equation}

Substituting (65) into (63) gives the final excited-state probability,

\begin{equation}
\boxed{
P_e(t)=\sin^2\left(\frac{\Omega_i t}{2}\right)
}.
\tag{66}
\end{equation}

The corresponding magnitude of the excited-state amplitude is

\begin{equation}
|\zeta_{e,i}(t)|
=
\left|\sin\left(\frac{\Omega_i t}{2}\right)\right|.
\tag{67}
\end{equation}

The oscillation argument in (66) is equivalently written as

\begin{equation}
\frac{\Omega_i t}{2}
=
\frac{|C_i|t}{2\hbar}.
\tag{68}
\end{equation}

Thus, the excited-state probability is completely determined by the magnitude of the effective interaction coefficient $C_i$ through the Rabi frequency $\Omega_i$.

\begin{equation}
0\leq P_e(t)\leq 1.
\tag{69}
\end{equation}

For the ground state, (42) and (58) give
\(\dot b_i=-\frac{jC_i^{*}}{2\hbar}a_i=-\frac{|C_i|}{2\hbar}\sin\!\left(\frac{|C_i|t}{2\hbar}\right)\).
With \(b_i(0)=1\), this integrates to \(b_i(t)=\cos\!\left(\frac{|C_i|t}{2\hbar}\right)\), so that
\(|\zeta_{g,i}(t)|^2=\cos^2\!\left(\frac{\Omega_i t}{2}\right)\).
The normalization condition is therefore satisfied because the corresponding ground-state probability obeys

\begin{equation}
|\zeta_{g,i}(t)|^2+P_e(t)=1.
\tag{70}
\end{equation}

\section{Expression for the Rabi Frequency}

From the definition of the effective interaction coefficient,

\begin{equation}
C_i
=
\sum_{k=1}^{K}
\sum_{l=1}^{L_k}
\boldsymbol{\mu}_{eg}^{T}
\boldsymbol{\epsilon}_{i,k,l}
\alpha_{l,k}s_k
 e^{-j(i-1)\phi_{l,k}}.
\tag{71}
\end{equation}

Using the definition of the Rabi frequency in (64),

\begin{equation}
\Omega_i
=
\frac{1}{\hbar}
\left|
\sum_{k=1}^{K}
\sum_{l=1}^{L_k}
\boldsymbol{\mu}_{eg}^{T}
\boldsymbol{\epsilon}_{i,k,l}
\alpha_{l,k}s_k
 e^{-j(i-1)\phi_{l,k}}
\right|.
\tag{72}
\end{equation}

Since $s_k$ does not depend on the path index $l$, it can be factored outside the inner summation:

\begin{equation}
\Omega_i
=
\frac{1}{\hbar}
\left|
\sum_{k=1}^{K}
 s_k
\sum_{l=1}^{L_k}
\boldsymbol{\mu}_{eg}^{T}
\boldsymbol{\epsilon}_{i,k,l}
\alpha_{l,k}
 e^{-j(i-1)\phi_{l,k}}
\right|.
\tag{73}
\end{equation}

Equivalently,

\begin{equation}
\Omega_i
=
\left|
\sum_{k=1}^{K}
 s_k
\left(
\sum_{l=1}^{L_k}
\frac{1}{\hbar}
\boldsymbol{\mu}_{eg}^{T}
\boldsymbol{\epsilon}_{i,k,l}
\alpha_{l,k}
 e^{-j(i-1)\phi_{l,k}}
\right)
\right|.
\tag{74}
\end{equation}

\begin{equation}
\boxed{
g_{i,k}
=
\sum_{l=1}^{L_k}
\frac{1}{\hbar}
\boldsymbol{\mu}_{eg}^{T}
\boldsymbol{\epsilon}_{i,k,l}
\alpha_{l,k}
e^{-j(i-1)\phi_{l,k}}
}.
\tag{75}
\end{equation}

Therefore,

\begin{equation}
\boxed{
\Omega_i
=
\left|
\sum_{k=1}^{K}
g_{i,k}s_k
\right|
}.
\tag{76}
\end{equation}

\section{Received Signal Model}

For the $p$-th time slot, the measured signal at vapor cell $i$ is modeled as

\begin{equation}
\boxed{
y_{i,p}
=
\left|
\sum_{k=1}^{K}
g_{i,k}s_{k,p}
+
b_{i,p}
+
n_{i,p}
\right|
}.
\tag{77}
\end{equation}

For a single communication user $k$, the contribution received at vapor cell $i$ during time slot $p$ is

\begin{equation}
g_{i,k}s_{k,p}.
\tag{78}
\end{equation}

Thus, the channel coefficient represents the complex coupling between the transmitted symbol and the received contribution:

\begin{equation}
\text{received contribution}
=
\text{channel coefficient}
\times
\text{transmitted symbol}.
\tag{79}
\end{equation}

For $K$ communication users, the contributions arriving at the same vapor cell are superposed linearly. Hence,

\begin{equation}
a_{i,p}
=
\sum_{k=1}^{K}
g_{i,k}s_{k,p},
\tag{80}
\end{equation}

where $a_{i,p}$ denotes the total desired wireless signal at vapor cell $i$ and time slot $p$.

The channel coefficient in (75) is obtained from the atomic coupling term,

\begin{equation}
\frac{1}{\hbar}
\boldsymbol{\mu}_{eg}^{T}
\boldsymbol{\epsilon}_{i,k,l}
\alpha_{l,k}
e^{-j(i-1)\phi_{l,k}},
\tag{81}
\end{equation}

which contains the atomic transition dipole coupling, polarization, path amplitude, and spatial propagation phase.

\section{Reference Signal}

In addition to the communication signals, a reference source is introduced.

Let the reference source transmit the symbol

\begin{equation}
s_{b,p},
\tag{82}
\end{equation}

where the subscript $b$ denotes the reference beam.

The reference source produces an electromagnetic field at vapor cell $i$. Its contribution is therefore written in the same form as a single propagation path:

\begin{equation}
b_{i,p}
=
g_{b,i}s_{b,p},
\tag{83}
\end{equation}

where $g_{b,i}$ represents the channel coupling from the reference source to vapor cell $i$.

For the communication users, multiple propagation paths are considered:

\begin{equation}
g_{i,k}
=
\sum_{l=1}^{L_k}
\left(
\text{contribution of path }l
\right).
\tag{84}
\end{equation}

For the reference source, the model uses one effective propagation path. Therefore, the reference channel contains one effective set of parameters,

\begin{equation}
\boldsymbol{\epsilon}_{b,i},
\qquad
\alpha_b,
\qquad
\phi_b.
\tag{85}
\end{equation}

Starting from the same single-path channel structure,

\begin{equation}
g_{i,k}^{(\mathrm{single})}
=
\frac{1}{\hbar}
\boldsymbol{\mu}_{eg}^{T}
\boldsymbol{\epsilon}_{i,k}
\alpha_k
e^{-j(i-1)\phi_k},
\tag{86}
\end{equation}

the reference channel is obtained by replacing the generic path parameters with the reference parameters:

\begin{equation}
\boldsymbol{\epsilon}_{i,k}
\rightarrow
\boldsymbol{\epsilon}_{b,i},
\qquad
\alpha_k
\rightarrow
\alpha_b,
\qquad
\phi_k
\rightarrow
\phi_b.
\tag{87}
\end{equation}

Hence,

\begin{equation}
\boxed{
g_{b,i}
=
\frac{1}{\hbar}
\boldsymbol{\mu}_{eg}^{T}
\boldsymbol{\epsilon}_{b,i}
\alpha_b
e^{-j(i-1)\phi_b}
}.
\tag{88}
\end{equation}

Substituting (88) into (83),

\begin{equation}
\boxed{
b_{i,p}
=
s_{b,p}
\left[
\frac{1}{\hbar}
\boldsymbol{\mu}_{eg}^{T}
\boldsymbol{\epsilon}_{b,i}
\alpha_b
e^{-j(i-1)\phi_b}
\right]
}.
\tag{89}
\end{equation}

\section{Complete Received Signal}

The receiver noise at vapor cell $i$ and time slot $p$ is denoted by

\begin{equation}
n_{i,p}
\sim
\mathcal{CN}(0,\sigma^2).
\tag{90}
\end{equation}

The total complex quantity experienced by the atomic receiver is therefore

\begin{equation}
a_{i,p}
+
b_{i,p}
+
n_{i,p},
\tag{91}
\end{equation}

or, using (80),

\begin{equation}
\sum_{k=1}^{K}
g_{i,k}s_{k,p}
+
b_{i,p}
+
n_{i,p}.
\tag{92}
\end{equation}

The Rydberg atomic receiver provides a magnitude-based measurement of this total coupling. Consequently,

\begin{equation}
\boxed{
y_{i,p}
=
\left|
a_{i,p}
+
b_{i,p}
+
n_{i,p}
\right|
}.
\tag{93}
\end{equation}

Using (80),

\begin{equation}
\boxed{
y_{i,p}
=
\left|
\sum_{k=1}^{K}
g_{i,k}s_{k,p}
+
b_{i,p}
+
n_{i,p}
\right|
}.
\tag{94}
\end{equation}

\section{Strong-Reference Approximation}

Define

\begin{equation}
x_{i,p}
=
a_{i,p}+n_{i,p}.
\tag{95}
\end{equation}

Then (93) becomes

\begin{equation}
y_{i,p}
=
|b_{i,p}+x_{i,p}|.
\tag{96}
\end{equation}

Squaring both sides,

\begin{equation}
y_{i,p}^{\,2}
=
|b_{i,p}+x_{i,p}|^2.
\tag{97}
\end{equation}

For any complex quantity $z$,

\begin{equation}
|z|^2=zz^*.
\tag{98}
\end{equation}

Therefore,

\begin{equation}
y_{i,p}^{\,2}
=
(b_{i,p}+x_{i,p})
(b_{i,p}+x_{i,p})^*.
\tag{99}
\end{equation}

Using

\begin{equation}
(b+x)^*=b^*+x^*,
\tag{100}
\end{equation}

gives

\begin{align}
y_{i,p}^{\,2}
&=
(b_{i,p}+x_{i,p})
(b_{i,p}^*+x_{i,p}^*)
\nonumber\\
&=
b_{i,p}b_{i,p}^*
+
b_{i,p}x_{i,p}^*
+
x_{i,p}b_{i,p}^*
+
x_{i,p}x_{i,p}^*.
\tag{101}
\end{align}

Since

\begin{equation}
|b_{i,p}|^2=b_{i,p}b_{i,p}^*,
\qquad
|x_{i,p}|^2=x_{i,p}x_{i,p}^*,
\tag{102}
\end{equation}

and

\begin{equation}
z+z^*=2\operatorname{Re}(z),
\tag{103}
\end{equation}

the cross terms satisfy

\begin{equation}
b_{i,p}x_{i,p}^*
+
b_{i,p}^*x_{i,p}
=
2\operatorname{Re}
\left(
b_{i,p}^*x_{i,p}
\right).
\tag{104}
\end{equation}

Hence,

\begin{equation}
\boxed{
y_{i,p}^{\,2}
=
|b_{i,p}|^2
+
|x_{i,p}|^2
+
2\operatorname{Re}
\left(
b_{i,p}^*x_{i,p}
\right)
}.
\tag{105}
\end{equation}

Taking the square root,

\begin{equation}
y_{i,p}
=
\sqrt{
|b_{i,p}|^2
+
|x_{i,p}|^2
+
2\operatorname{Re}
\left(
b_{i,p}^*x_{i,p}
\right)
}.
\tag{106}
\end{equation}

Factor $|b_{i,p}|^2$ from the square root:

\begin{equation}
y_{i,p}
=
|b_{i,p}|
\sqrt{
1+
\frac{|x_{i,p}|^2}{|b_{i,p}|^2}
+
\frac{
2\operatorname{Re}(b_{i,p}^*x_{i,p})
}{
|b_{i,p}|^2
}
}.
\tag{107}
\end{equation}

Using

\begin{equation}
|b_{i,p}|^2=b_{i,p}b_{i,p}^*,
\tag{108}
\end{equation}

the last term becomes

\begin{equation}
\frac{
b_{i,p}^*x_{i,p}
}{
b_{i,p}b_{i,p}^*
}
=
\frac{x_{i,p}}{b_{i,p}}.
\tag{109}
\end{equation}

Thus,

\begin{equation}
\boxed{
y_{i,p}
=
|b_{i,p}|
\sqrt{
1+
\frac{|x_{i,p}|^2}{|b_{i,p}|^2}
+
2\operatorname{Re}
\left(
\frac{x_{i,p}}{b_{i,p}}
\right)
}
}.
\tag{110}
\end{equation}

\section{Taylor Expansion Under Strong Reference}

The strong-reference assumption is

\begin{equation}
|b_{i,p}|
\gg
|x_{i,p}|
=
|a_{i,p}+n_{i,p}|.
\tag{111}
\end{equation}

Consequently,

\begin{equation}
\left|
\frac{x_{i,p}}{b_{i,p}}
\right|
\ll 1.
\tag{112}
\end{equation}

Define

\begin{equation}
r_{i,p}
=
\frac{x_{i,p}}{b_{i,p}}.
\tag{113}
\end{equation}

Then

\begin{equation}
|r_{i,p}|\ll1.
\tag{114}
\end{equation}

Equation (110) becomes

\begin{equation}
y_{i,p}
=
|b_{i,p}|
\sqrt{
1+
|r_{i,p}|^2
+
2\operatorname{Re}(r_{i,p})
}.
\tag{115}
\end{equation}

For $|u|\ll1$, the first-order Taylor expansion is

\begin{equation}
\sqrt{1+u}
\approx
1+\frac{u}{2}.
\tag{116}
\end{equation}

Let

\begin{equation}
u
=
|r_{i,p}|^2
+
2\operatorname{Re}(r_{i,p}).
\tag{117}
\end{equation}

Then,

\begin{equation}
y_{i,p}
\approx
|b_{i,p}|
\left[
1+
\frac{1}{2}
\left(
|r_{i,p}|^2
+
2\operatorname{Re}(r_{i,p})
\right)
\right].
\tag{118}
\end{equation}

Therefore,

\begin{equation}
y_{i,p}
\approx
|b_{i,p}|
\left[
1+
\frac{1}{2}|r_{i,p}|^2
+
\operatorname{Re}(r_{i,p})
\right].
\tag{119}
\end{equation}

Substituting (113),

\begin{equation}
y_{i,p}
\approx
|b_{i,p}|
\left[
1+
\frac{1}{2}
\frac{|x_{i,p}|^2}{|b_{i,p}|^2}
+
\operatorname{Re}
\left(
\frac{x_{i,p}}{b_{i,p}}
\right)
\right].
\tag{120}
\end{equation}

Expanding the factor $|b_{i,p}|$,

\begin{equation}
y_{i,p}
\approx
|b_{i,p}|
+
\frac{|x_{i,p}|^2}{2|b_{i,p}|}
+
|b_{i,p}|
\operatorname{Re}
\left(
\frac{x_{i,p}}{b_{i,p}}
\right).
\tag{121}
\end{equation}

The second term is second order in the small quantity $x_{i,p}/b_{i,p}$, whereas the third term is first order. Under the strong-reference approximation, the second-order term is neglected.
(A complete second-order expansion of \(\sqrt{1+u}\) would also contain
\(-\tfrac12|b_{i,p}|\operatorname{Re}^2(x_{i,p}/b_{i,p})\); it is neglected for the same reason.)

Hence,

\begin{equation}
\boxed{
y_{i,p}
\approx
|b_{i,p}|
+
|b_{i,p}|
\operatorname{Re}
\left(
\frac{x_{i,p}}{b_{i,p}}
\right)
}.
\tag{122}
\end{equation}

Since

\begin{equation}
x_{i,p}
=
a_{i,p}+n_{i,p},
\tag{123}
\end{equation}

we obtain

\begin{equation}
y_{i,p}
\approx
|b_{i,p}|
+
|b_{i,p}|
\operatorname{Re}
\left(
\frac{a_{i,p}+n_{i,p}}{b_{i,p}}
\right).
\tag{124}
\end{equation}

\section{Reference Magnitude and Phase Representation}

Represent the reference signal in magnitude-phase form as

\begin{equation}
b_{i,p}
=
|b_{i,p}|e^{jz_{i,p}},
\tag{125}
\end{equation}

where

\begin{equation}
z_{i,p}
=
\angle b_{i,p}
=
\arg(b_{i,p}).
\tag{126}
\end{equation}

Therefore,

\begin{equation}
\frac{|b_{i,p}|}{b_{i,p}}
=
e^{-jz_{i,p}}.
\tag{127}
\end{equation}

Using (127) in (124),

\begin{equation}
y_{i,p}
\approx
|b_{i,p}|
+
\operatorname{Re}
\left(
e^{-jz_{i,p}}
(a_{i,p}+n_{i,p})
\right).
\tag{128}
\end{equation}

Using the linearity of the real-part operator,

\begin{equation}
\operatorname{Re}(A+B)
=
\operatorname{Re}(A)
+
\operatorname{Re}(B),
\tag{129}
\end{equation}

gives

\begin{equation}
y_{i,p}
\approx
|b_{i,p}|
+
\operatorname{Re}
\left(
e^{-jz_{i,p}}a_{i,p}
\right)
+
\operatorname{Re}
\left(
e^{-jz_{i,p}}n_{i,p}
\right).
\tag{130}
\end{equation}

Define the effective real-valued noise term

\begin{equation}
\bar n_{i,p}
=
\operatorname{Re}
\left(
e^{-jz_{i,p}}n_{i,p}
\right).
\tag{131}
\end{equation}

Since $n_{i,p}$ is circularly symmetric, $e^{-jz_{i,p}}n_{i,p}\sim\mathcal{CN}(0,\sigma^2)$,
and its real part carries half of the variance:

\begin{equation}
\boxed{
\bar n_{i,p}
\sim
\mathcal{N}\!\left(0,\frac{\sigma^2}{2}\right)
}.
\tag{131a}
\end{equation}

Thus,

\begin{equation}
\boxed{
y_{i,p}
\approx
\operatorname{Re}
\left(
e^{-jz_{i,p}}a_{i,p}
\right)
+
|b_{i,p}|
+
\bar n_{i,p}
}.
\tag{132}
\end{equation}

Finally, using

\begin{equation}
a_{i,p}
=
\sum_{k=1}^{K}
g_{i,k}s_{k,p},
\tag{133}
\end{equation}

the linearized measurement model becomes

\begin{equation}
\boxed{
y_{i,p}
\approx
\operatorname{Re}
\left(
e^{-jz_{i,p}}
\sum_{k=1}^{K}
g_{i,k}s_{k,p}
\right)
+
|b_{i,p}|
+
\bar n_{i,p}
}.
\tag{134}
\end{equation}

% ============================================================
% END OF PAGES 21--30
% ===========================================================

% ============================================================
% PAGES 31--40
% ============================================================

\section{Matrix Representation of the Linearized Measurement Model}

From the scalar linearized measurement model obtained previously,

\begin{equation}
y_{i,p}
\approx
\operatorname{Re}
\left(
e^{-jz_{i,p}}
\sum_{k=1}^{K}
g_{i,k}s_{k,p}
\right)
+
|b_{i,p}|
+
\bar n_{i,p}.
\tag{135}
\end{equation}

The indices are defined as

\[
i=1,2,\ldots,I
\]

for the $I$ vapor cells or receiving antennas,

\[
k=1,2,\ldots,K
\]

for the $K$ users, and

\[
p=1,2,\ldots,P
\]

for the $P$ transmitted symbols or time slots.

Thus, $y_{i,p}$ represents one scalar magnitude measurement corresponding to vapor cell $i$ and time slot $p$.

Collecting all measurements gives the matrix

\begin{equation}
\boxed{
Y=[y_{i,p}]
\in
\mathbb{R}^{I\times P}
}.
\tag{136}
\end{equation}

Explicitly,

\begin{equation}
Y=
\begin{bmatrix}
y_{1,1} & y_{1,2} & \cdots & y_{1,P}\\
y_{2,1} & y_{2,2} & \cdots & y_{2,P}\\
\vdots & \vdots & \ddots & \vdots\\
y_{I,1} & y_{I,2} & \cdots & y_{I,P}
\end{bmatrix}.
\tag{137}
\end{equation}

Since $y_{i,p}$ is a magnitude measurement,

\begin{equation}
y_{i,p}\in\mathbb{R}.
\tag{138}
\end{equation}

\subsection{Construction of the Phase Matrix}

The scalar model contains the phase factor

\begin{equation}
e^{-jz_{i,p}}.
\tag{139}
\end{equation}

Define the phase matrix

\begin{equation}
\boxed{
Z=
[e^{-jz_{i,p}}]
\in
\mathbb{C}^{I\times P}
}.
\tag{140}
\end{equation}

The matrix is therefore

\begin{equation}
Z=
\begin{bmatrix}
e^{-jz_{1,1}} & e^{-jz_{1,2}} & \cdots & e^{-jz_{1,P}}\\
e^{-jz_{2,1}} & e^{-jz_{2,2}} & \cdots & e^{-jz_{2,P}}\\
\vdots & \vdots & \ddots & \vdots\\
e^{-jz_{I,1}} & e^{-jz_{I,2}} & \cdots & e^{-jz_{I,P}}
\end{bmatrix}.
\tag{141}
\end{equation}

The phase variable is determined from the reference signal according to

\begin{equation}
z_{i,p}
=
\angle b_{i,p}.
\tag{142}
\end{equation}

\subsection{Construction of the Channel Matrix}

The channel coefficient between vapor cell $i$ and user $k$ is $g_{i,k}$.

There are $I$ receiving cells and $K$ users. Hence, the channel matrix is

\begin{equation}
\boxed{
G=[g_{i,k}]
\in
\mathbb{C}^{I\times K}
}.
\tag{143}
\end{equation}

Explicitly,

\begin{equation}
G=
\begin{bmatrix}
g_{1,1} & g_{1,2} & \cdots & g_{1,K}\\
g_{2,1} & g_{2,2} & \cdots & g_{2,K}\\
\vdots & \vdots & \ddots & \vdots\\
g_{I,1} & g_{I,2} & \cdots & g_{I,K}
\end{bmatrix}.
\tag{144}
\end{equation}

Similarly, the transmitted-symbol matrix is defined as

\begin{equation}
\boxed{
S=[s_{k,p}]
\in
\mathbb{C}^{K\times P}
}.
\tag{145}
\end{equation}

Explicitly,

\begin{equation}
S=
\begin{bmatrix}
s_{1,1} & s_{1,2} & \cdots & s_{1,P}\\
s_{2,1} & s_{2,2} & \cdots & s_{2,P}\\
\vdots & \vdots & \ddots & \vdots\\
s_{K,1} & s_{K,2} & \cdots & s_{K,P}
\end{bmatrix}.
\tag{146}
\end{equation}

The matrix product $GS$ has dimensions

\begin{equation}
G S:
\qquad
(I\times K)(K\times P)
=
I\times P.
\tag{147}
\end{equation}

Its $(i,p)$-th element is

\begin{equation}
(GS)_{i,p}
=
\sum_{k=1}^{K}
g_{i,k}s_{k,p}.
\tag{148}
\end{equation}

Therefore,

\begin{equation}
(GS)_{i,p}=a_{i,p}.
\tag{149}
\end{equation}

\subsection{Hadamard Product with the Phase Matrix}

The phase factor $e^{-jz_{i,p}}$ has to multiply the corresponding $(i,p)$-th element of $GS$.

For two matrices $A$ and $B$ having the same dimensions, the Hadamard product is defined element-wise as

\begin{equation}
(A\circ B)_{i,p}
=
A_{i,p}B_{i,p}.
\tag{150}
\end{equation}

Therefore,

\begin{equation}
[Z\circ(GS)]_{i,p}
=
Z_{i,p}(GS)_{i,p}.
\tag{151}
\end{equation}

Using

\begin{equation}
Z_{i,p}=e^{-jz_{i,p}},
\tag{152}
\end{equation}

and (148),

\begin{equation}
[Z\circ(GS)]_{i,p}
=
e^{-jz_{i,p}}
\sum_{k=1}^{K}
g_{i,k}s_{k,p}.
\tag{153}
\end{equation}

Applying the real-part operator gives

\begin{equation}
\left(
\operatorname{Re}[Z\circ(GS)]
\right)_{i,p}
=
\operatorname{Re}
\left(
e^{-jz_{i,p}}
\sum_{k=1}^{K}
g_{i,k}s_{k,p}
\right).
\tag{154}
\end{equation}

\subsection{Construction of the Reference Magnitude Matrix}

The scalar measurement model contains the reference magnitude

\begin{equation}
|b_{i,p}|.
\tag{155}
\end{equation}

Define the reference matrix as

\begin{equation}
\boxed{
B=[b_{i,p}]
\in
\mathbb{C}^{I\times P}
}.
\tag{156}
\end{equation}

The corresponding element-wise absolute-value matrix is

\begin{equation}
\boxed{
|B|
=
[|b_{i,p}|]
\in
\mathbb{R}^{I\times P}
}.
\tag{157}
\end{equation}

Here $|B|$ denotes the element-wise absolute value and not a matrix norm. Thus,

\begin{equation}
|B|_{i,p}
=
|b_{i,p}|.
\tag{158}
\end{equation}

\subsection{Construction of the Noise Matrix}

The effective real-valued noise term is $\bar n_{i,p}$. Define

\begin{equation}
\boxed{
N=[\bar n_{i,p}]
\in
\mathbb{R}^{I\times P}
}.
\tag{159}
\end{equation}

Explicitly,

\begin{equation}
N=
\begin{bmatrix}
\bar n_{1,1} & \bar n_{1,2} & \cdots & \bar n_{1,P}\\
\bar n_{2,1} & \bar n_{2,2} & \cdots & \bar n_{2,P}\\
\vdots & \vdots & \ddots & \vdots\\
\bar n_{I,1} & \bar n_{I,2} & \cdots & \bar n_{I,P}
\end{bmatrix}.
\tag{160}
\end{equation}

\section{Complete Matrix Measurement Model}

The scalar linearized model is

\begin{equation}
y_{i,p}
\approx
\operatorname{Re}
\left(
e^{-jz_{i,p}}
\sum_{k=1}^{K}
g_{i,k}s_{k,p}
\right)
+
|b_{i,p}|
+
\bar n_{i,p}.
\tag{161}
\end{equation}

Using the definitions above, the element-wise representation becomes

\begin{equation}
y_{i,p}
=
\operatorname{Re}
\left(
[Z\circ(GS)]_{i,p}
\right)
+
|B|_{i,p}
+
N_{i,p}.
\tag{162}
\end{equation}

Collecting all elements into matrix form gives

\begin{equation}
\boxed{
Y
=
\operatorname{Re}
[Z\circ(GS)]
+
|B|
+
N
}.
\tag{163}
\end{equation}

The Hadamard product is commutative,

\begin{equation}
Z\circ(GS)
=
(GS)\circ Z.
\tag{164}
\end{equation}

Therefore, either ordering produces the same element-wise product.

The use of the Hadamard product is required because each phase coefficient $e^{-jz_{i,p}}$ multiplies the corresponding measurement position $(i,p)$ rather than performing an ordinary matrix multiplication.

\section{Channel Estimation Problem}

The matrix measurement model is

\begin{equation}
Y
=
\operatorname{Re}
[Z\circ(GS)]
+
|B|
+
N.
\tag{165}
\end{equation}

The objective is to estimate the unknown channel matrix

\begin{equation}
G=[g_{i,k}]
\tag{166}
\end{equation}

from the available measurements.

The known quantities are

\begin{equation}
Y,\qquad Z,\qquad B,\qquad S,
\tag{167}
\end{equation}

while the channel matrix $G$ is unknown.

Let $G^{t}$ denote the current estimate of the channel matrix at iteration $t$.

Using this estimate, the predicted measurement is

\begin{equation}
\boxed{
\widehat{Y}(G^t)
=
\operatorname{Re}
[Z\circ(G^tS)]
+
|B|
}.
\tag{168}
\end{equation}

The prediction error is defined as

\begin{equation}
\boxed{
E(G^t)
=
Y-\widehat{Y}(G^t)
}.
\tag{169}
\end{equation}

Substituting (168),

\begin{equation}
\boxed{
E(G^t)
=
Y-|B|
-
\operatorname{Re}
[Z\circ(G^tS)]
}.
\tag{170}
\end{equation}

The dimensions are

\begin{equation}
E(G^t)\in\mathbb{R}^{I\times P}.
\tag{171}
\end{equation}

\section{Least-Squares Objective}

The objective is to obtain a channel estimate whose predicted measurements are as close as possible to the actual measurements.

The squared Frobenius norm is used to measure the total squared error:

\begin{equation}
\boxed{
\mathcal{L}(G)
=
\|E(G)\|_F^2
}.
\tag{172}
\end{equation}

Using (170),

\begin{equation}
\boxed{
\mathcal{L}(G)
=
\left\|
Y-|B|
-
\operatorname{Re}
[Z\circ(GS)]
\right\|_F^2
}.
\tag{173}
\end{equation}

For an $I\times P$ matrix,

\begin{equation}
\|E\|_F^2
=
\sum_{i=1}^{I}
\sum_{p=1}^{P}
|E_{i,p}|^2.
\tag{174}
\end{equation}

Since $E$ is real-valued,

\begin{equation}
\|E\|_F^2
=
\sum_{i=1}^{I}
\sum_{p=1}^{P}
E_{i,p}^2.
\tag{175}
\end{equation}

Therefore, the channel estimate is obtained by solving

\begin{equation}
\boxed{
\widehat{G}
=
\arg\min_G
\mathcal{L}(G)
}.
\tag{176}
\end{equation}

Equivalently,

\begin{equation}
\boxed{
\widehat{G}
=
\arg\min_G
\left\|
Y-|B|
-
\operatorname{Re}
[Z\circ(GS)]
\right\|_F^2
}.
\tag{177}
\end{equation}

\section{Element-by-Element Error Expression}

From (170), the $(i,p)$-th error element is

\begin{equation}
E_{i,p}
=
y_{i,p}
-
|b_{i,p}|
-
\operatorname{Re}
\left[
Z_{i,p}
\sum_{r=1}^{K}
g_{i,r}s_{r,p}
\right].
\tag{178}
\end{equation}

Here $r$ is a dummy summation index.

The channel coefficient with respect to which differentiation will be performed is $g_{i,k}$. The index $k$ is therefore kept fixed during the differentiation.

The cost function can be written element-wise as

\begin{equation}
\boxed{
\mathcal{L}
=
\sum_{i=1}^{I}
\sum_{p=1}^{P}
E_{i,p}^2
}.
\tag{179}
\end{equation}

To obtain the gradient with respect to the complex channel matrix, differentiation is performed with respect to the conjugate variable $g_{i,k}^{*}$.

\section{Wirtinger Differentiation}

Since $G$ is complex-valued, the loss depends on both $G$ and $G^*$. Therefore, Wirtinger differentiation is used.

For the selected channel coefficient $g_{i,k}$,

\begin{equation}
\frac{\partial\mathcal{L}}
{\partial g_{i,k}^{*}}
=
\sum_{p=1}^{P}
2E_{i,p}
\frac{\partial E_{i,p}}
{\partial g_{i,k}^{*}}.
\tag{180}
\end{equation}

This follows from the chain rule applied to

\begin{equation}
\mathcal{L}
=
\sum_{i=1}^{I}
\sum_{p=1}^{P}
E_{i,p}^{2}.
\tag{181}
\end{equation}

For the selected $g_{i,k}$, only the errors corresponding to the same receiving-cell index $i$ depend on $g_{i,k}$.

From (178),

\begin{equation}
E_{i,p}
=
y_{i,p}
-
|b_{i,p}|
-
\operatorname{Re}
\left[
Z_{i,p}
\sum_{r=1}^{K}
g_{i,r}s_{r,p}
\right].
\tag{182}
\end{equation}

Using

\begin{equation}
\operatorname{Re}(x)
=
\frac{x+x^*}{2},
\tag{183}
\end{equation}

we obtain

\begin{align}
E_{i,p}
&=
y_{i,p}
-
|b_{i,p}|
\nonumber\\
&\quad
-
\frac{1}{2}
\left[
Z_{i,p}
\sum_{r=1}^{K}
g_{i,r}s_{r,p}
+
Z_{i,p}^{*}
\sum_{r=1}^{K}
g_{i,r}^{*}s_{r,p}^{*}
\right].
\tag{184}
\end{align}

Differentiate (184) with respect to $g_{i,k}^{*}$:

\begin{equation}
\frac{\partial E_{i,p}}
{\partial g_{i,k}^{*}}
=
-\frac{1}{2}
\frac{\partial}
{\partial g_{i,k}^{*}}
\left[
Z_{i,p}
\sum_{r=1}^{K}
g_{i,r}s_{r,p}
+
Z_{i,p}^{*}
\sum_{r=1}^{K}
g_{i,r}^{*}s_{r,p}^{*}
\right].
\tag{185}
\end{equation}

The first term contains $g_{i,r}$ rather than $g_{i,r}^{*}$. Under Wirtinger differentiation,

\begin{equation}
\frac{\partial g_{i,r}}
{\partial g_{i,k}^{*}}
=
0.
\tag{186}
\end{equation}

Therefore,

\begin{equation}
\frac{\partial}
{\partial g_{i,k}^{*}}
\left(
Z_{i,p}
\sum_{r=1}^{K}
g_{i,r}s_{r,p}
\right)
=
0.
\tag{187}
\end{equation}

For the second term,

\begin{equation}
\frac{\partial}
{\partial g_{i,k}^{*}}
\left(
Z_{i,p}^{*}
\sum_{r=1}^{K}
g_{i,r}^{*}s_{r,p}^{*}
\right)
=
Z_{i,p}^{*}
\sum_{r=1}^{K}
\frac{\partial g_{i,r}^{*}}
{\partial g_{i,k}^{*}}
s_{r,p}^{*}.
\tag{188}
\end{equation}

The derivative satisfies

\begin{equation}
\frac{\partial g_{i,r}^{*}}
{\partial g_{i,k}^{*}}
=
\begin{cases}
1, & r=k,\\
0, & r\neq k.
\end{cases}
\tag{189}
\end{equation}

Consequently,

\begin{equation}
\frac{\partial}
{\partial g_{i,k}^{*}}
\left(
Z_{i,p}^{*}
\sum_{r=1}^{K}
g_{i,r}^{*}s_{r,p}^{*}
\right)
=
Z_{i,p}^{*}s_{k,p}^{*}.
\tag{190}
\end{equation}

Substituting (187) and (190) into (185),

\begin{equation}
\boxed{
\frac{\partial E_{i,p}}
{\partial g_{i,k}^{*}}
=
-\frac{1}{2}
Z_{i,p}^{*}s_{k,p}^{*}
}.
\tag{191}
\end{equation}

\section{Derivative of the Least-Squares Cost}

Substituting (191) into (180),

\begin{align}
\frac{\partial\mathcal{L}}
{\partial g_{i,k}^{*}}
&=
\sum_{p=1}^{P}
2E_{i,p}
\left(
-\frac{1}{2}
Z_{i,p}^{*}s_{k,p}^{*}
\right)
\nonumber\\
&=
-\sum_{p=1}^{P}
E_{i,p}
Z_{i,p}^{*}
s_{k,p}^{*}.
\tag{192}
\end{align}

Hence,

\begin{equation}
\boxed{
\frac{\partial\mathcal{L}}
{\partial g_{i,k}^{*}}
=
-\sum_{p=1}^{P}
E_{i,p}
Z_{i,p}^{*}
s_{k,p}^{*}
}.
\tag{193}
\end{equation}

The above operation is performed for every channel coefficient

\begin{equation}
g_{i,k},
\qquad
i=1,\ldots,I,
\qquad
k=1,\ldots,K.
\tag{194}
\end{equation}

Thus, all derivatives are collected into the gradient matrix

\begin{equation}
\boxed{
\frac{\partial\mathcal{L}}
{\partial G^{*}}
=
\begin{bmatrix}
\dfrac{\partial\mathcal{L}}{\partial g_{1,1}^{*}}
&
\dfrac{\partial\mathcal{L}}{\partial g_{1,2}^{*}}
&
\cdots
&
\dfrac{\partial\mathcal{L}}{\partial g_{1,K}^{*}}
\\
\dfrac{\partial\mathcal{L}}{\partial g_{2,1}^{*}}
&
\dfrac{\partial\mathcal{L}}{\partial g_{2,2}^{*}}
&
\cdots
&
\dfrac{\partial\mathcal{L}}{\partial g_{2,K}^{*}}
\\
\vdots
&
\vdots
&
\ddots
&
\vdots
\\
\dfrac{\partial\mathcal{L}}{\partial g_{I,1}^{*}}
&
\dfrac{\partial\mathcal{L}}{\partial g_{I,2}^{*}}
&
\cdots
&
\dfrac{\partial\mathcal{L}}{\partial g_{I,K}^{*}}
\end{bmatrix}
\in
\mathbb{C}^{I\times K}
}.
\tag{195}
\end{equation}

\section{Conversion of the Gradient to Matrix Form}

From (193),

\begin{equation}
\frac{\partial\mathcal{L}}
{\partial g_{i,k}^{*}}
=
-\sum_{p=1}^{P}
E_{i,p}
Z_{i,p}^{*}
s_{k,p}^{*}.
\tag{196}
\end{equation}

Consider the matrix product

\begin{equation}
(E\circ Z^{*})S^{H}.
\tag{197}
\end{equation}

The dimensions are

\begin{equation}
E\in\mathbb{R}^{I\times P},
\qquad
Z^{*}\in\mathbb{C}^{I\times P},
\tag{198}
\end{equation}

therefore,

\begin{equation}
E\circ Z^{*}
\in
\mathbb{C}^{I\times P}.
\tag{199}
\end{equation}

Also,

\begin{equation}
S\in\mathbb{C}^{K\times P},
\qquad
S^{H}\in\mathbb{C}^{P\times K}.
\tag{200}
\end{equation}

Hence,

\begin{equation}
(E\circ Z^{*})S^{H}
:
(I\times P)(P\times K)
=
I\times K.
\tag{201}
\end{equation}

Consider the $(i,k)$-th element:

\begin{align}
[(E\circ Z^{*})S^{H}]_{i,k}
&=
\sum_{p=1}^{P}
(E\circ Z^{*})_{i,p}
(S^{H})_{p,k}.
\tag{202}
\end{align}

By the definition of the Hadamard product,

\begin{equation}
(E\circ Z^{*})_{i,p}
=
E_{i,p}Z_{i,p}^{*},
\tag{203}
\end{equation}

and

\begin{equation}
(S^{H})_{p,k}
=
S_{k,p}^{*}
=
s_{k,p}^{*}.
\tag{204}
\end{equation}

Therefore,

\begin{align}
[(E\circ Z^{*})S^{H}]_{i,k}
&=
\sum_{p=1}^{P}
E_{i,p}
Z_{i,p}^{*}
s_{k,p}^{*}.
\tag{205}
\end{align}

Comparing (205) with (193),

\begin{equation}
\boxed{
\frac{\partial\mathcal{L}}
{\partial G^{*}}
=
-(E\circ Z^{*})S^{H}
}.
\tag{206}
\end{equation}

\section{Gradient-Descent Channel Update}

The gradient-descent update rule is

\begin{equation}
\boxed{
G^{t+1}
=
G^{t}
-
\eta_t
\left.
\frac{\partial\mathcal{L}}
{\partial G^{*}}
\right|_{G=G^{t}}
}
\tag{207}
\end{equation}

where $\eta_t$ denotes the step size at iteration $t$.

Using (206),

\begin{align}
G^{t+1}
&=
G^{t}
-
\eta_t
\left[
-(E^{t}\circ Z^{*})S^{H}
\right]
\nonumber\\
&=
G^{t}
+
\eta_t
(E^{t}\circ Z^{*})S^{H}.
\tag{208}
\end{align}

The error at iteration $t$ is

\begin{equation}
\boxed{
E^{t}
=
Y-|B|
-
\operatorname{Re}
[Z\circ(G^{t}S)]
}.
\tag{209}
\end{equation}

Therefore, the complete channel update is

\begin{equation}
\boxed{
G^{t+1}
=
G^{t}
+
\eta_t
\left(
\left[
Y-|B|
-
\operatorname{Re}
[Z\circ(G^{t}S)]
\right]
\circ
Z^{*}
\right)
S^{H}
}.
\tag{210}
\end{equation}

The dimensions of the update are

\begin{equation}
G^{t}\in\mathbb{C}^{I\times K},
\tag{211}
\end{equation}

\begin{equation}
G^{t}S
\in
\mathbb{C}^{I\times P},
\tag{212}
\end{equation}

\begin{equation}
E^{t}
\in
\mathbb{R}^{I\times P},
\tag{213}
\end{equation}

\begin{equation}
E^{t}\circ Z^{*}
\in
\mathbb{C}^{I\times P},
\tag{214}
\end{equation}

and

\begin{equation}
(E^{t}\circ Z^{*})S^{H}
\in
\mathbb{C}^{I\times K}.
\tag{215}
\end{equation}

Thus, the gradient has the same dimensions as the channel matrix $G$.

\section{Gradient-Descent Channel Estimation Algorithm}

The gradient-descent channel estimation procedure is summarized as follows.

\begin{equation}
\boxed{
\text{Input: }
Y,\;S,\;Z,\;B,\;G^{0},\;\eta_t,\;\epsilon,\;T_{\max}
}
\tag{216}
\end{equation}

Initialize

\begin{equation}
t=0.
\tag{217}
\end{equation}

At iteration $t$, calculate the predicted measurement:

\begin{equation}
\boxed{
\widehat{Y}^{t}
=
\operatorname{Re}
[Z\circ(G^{t}S)]
+
|B|
}.
\tag{218}
\end{equation}

Calculate the measurement error:

\begin{equation}
\boxed{
E^{t}
=
Y-\widehat{Y}^{t}
}.
\tag{219}
\end{equation}

Equivalently,

\begin{equation}
\boxed{
E^{t}
=
Y-|B|
-
\operatorname{Re}
[Z\circ(G^{t}S)]
}.
\tag{220}
\end{equation}

Calculate the gradient:

\begin{equation}
\boxed{
D^{t}
=
-
(E^{t}\circ Z^{*})S^{H}
}.
\tag{221}
\end{equation}

Update the channel estimate:

\begin{equation}
\boxed{
G^{t+1}
=
G^{t}
-
\eta_tD^{t}
}.
\tag{222}
\end{equation}

Using (221),

\begin{equation}
\boxed{
G^{t+1}
=
G^{t}
+
\eta_t
(E^{t}\circ Z^{*})S^{H}
}.
\tag{223}
\end{equation}

The iteration index is then increased according to

\begin{equation}
t\leftarrow t+1.
\tag{224}
\end{equation}

The iterations continue until

\begin{equation}
\boxed{
\left\|
G^{t}-G^{t-1}
\right\|_{F}^{2}
<
\epsilon
}
\tag{225}
\end{equation}

or

\begin{equation}
t\geq T_{\max}.
\tag{226}
\end{equation}

The final channel estimate is

\begin{equation}
\boxed{
\widehat{G}=G^{t}
}.
\tag{227}
\end{equation}

\section{Algorithmic Form}

The complete gradient-descent procedure is

\begin{equation}
\boxed{
\begin{array}{l}
\textbf{Input: }
Y,\;S,\;Z,\;B,\;G^{0},\;\eta_t,\;\epsilon,\;T_{\max}
\\[2mm]
t=0
\\[2mm]
\textbf{Repeat:}
\\
\qquad
\widehat{Y}^{t}
=
\operatorname{Re}
[Z\circ(G^{t}S)]
+
|B|
\\[1mm]
\qquad
E^{t}
=
Y-\widehat{Y}^{t}
\\[1mm]
\qquad
D^{t}
=
-(E^{t}\circ Z^{*})S^{H}
\\[1mm]
\qquad
G^{t+1}
=
G^{t}
-
\eta_tD^{t}
\\[1mm]
\qquad
t\leftarrow t+1
\\[2mm]
\textbf{Until: }
\|G^{t}-G^{t-1}\|_{F}^{2}<\epsilon
\quad
\text{or}
\quad
t\geq T_{\max}
\\[2mm]
\textbf{Output: }
\widehat{G}=G^{t}.
\end{array}
}
\tag{228}
\end{equation}

% ============================================================
% END OF PAGES 31--40
% ===========================================================
% ============================================================
% PAGES 41--51
% ============================================================

\section{Biased Gerchberg--Saxton Algorithm for the 1D Antenna Array}
\label{sec:gs}

\subsection{1D Channel Estimation Model}

Consider an array consisting of $I$ atomic vapor cells and $K$ users.

For the $i$-th vapor cell, define the channel vector as

\begin{equation}
\boxed{
\mathbf g_i
=
\begin{bmatrix}
g_{i,1}\\
g_{i,2}\\
\vdots\\
g_{i,K}
\end{bmatrix}
\in
\mathbb C^{K\times 1}
}.
\tag{229}
\end{equation}

The known pilot matrix is

\begin{equation}
\boxed{
S=
\begin{bmatrix}
s_{1,1} & s_{1,2} & \cdots & s_{1,P}\\
s_{2,1} & s_{2,2} & \cdots & s_{2,P}\\
\vdots & \vdots & \ddots & \vdots\\
s_{K,1} & s_{K,2} & \cdots & s_{K,P}
\end{bmatrix}
\in
\mathbb C^{K\times P}
}.
\tag{230}
\end{equation}

The reference vector and noise vector are

\begin{equation}
\mathbf b_i\in\mathbb C^{P\times1},
\qquad
\mathbf w_i\in\mathbb C^{P\times1}.
\tag{231}
\end{equation}

The measured magnitude vector is

\begin{equation}
\boxed{
\mathbf y_i
=
\left|
S^T\mathbf g_i+\mathbf b_i+\mathbf w_i
\right|
\in
\mathbb R^{P\times1}
}.
\tag{232}
\end{equation}

The $p$-th element of $S^T\mathbf g_i$ is

\begin{equation}
\left[S^T\mathbf g_i\right]_p
=
\sum_{k=1}^{K}
s_{k,p}g_{i,k}.
\tag{233}
\end{equation}

The difficulty is that only the magnitude of the received complex signal is measured.

Define the unknown complex received vector

\begin{equation}
\boxed{
\mathbf x_i
=
S^T\mathbf g_i+\mathbf b_i+\mathbf w_i
}.
\tag{234}
\end{equation}

Then,

\begin{equation}
\mathbf y_i=|\mathbf x_i|.
\tag{235}
\end{equation}

If the phase of $\mathbf x_i$ were known, the complex received signal could be reconstructed as

\begin{equation}
\mathbf x_i
=
\mathbf y_i\circ e^{j\boldsymbol{\theta}_i},
\tag{236}
\end{equation}

where $\boldsymbol{\theta}_i$ denotes the unknown phase vector.

The GS procedure therefore alternates between phase estimation, complex-signal reconstruction, and channel estimation.

\section{Spectral Initialization}

The iterative GS algorithm depends on its initial channel estimate. A spectral initialization is therefore used to obtain an initial estimate of the channel.

The biased GS initialization constructs an augmented observation model, forms a weighted covariance matrix, extracts its principal eigenvector, estimates the corresponding magnitude, and removes the arbitrary global phase.

\subsection{Conversion to the Standard Spectral Form}

The channel measurement model is

\begin{equation}
\mathbf y_i
=
\left|
S^T\mathbf g_i+\mathbf b_i+\mathbf w_i
\right|.
\tag{237}
\end{equation}

The standard biased phase-retrieval model is written as

\begin{equation}
\mathbf y
=
\left|
A^H\mathbf g+\mathbf b+\mathbf w
\right|.
\tag{238}
\end{equation}

Comparing (237) and (238),

\begin{equation}
\boxed{
A^H=S^T
}.
\tag{239}
\end{equation}

Taking the Hermitian relation,

\begin{equation}
A=S^*.
\tag{240}
\end{equation}

Therefore, the channel model can be written as

\begin{equation}
\boxed{
\mathbf y_i
=
\left|
A^H\mathbf g_i+\mathbf b_i+\mathbf w_i
\right|
,
\qquad
A=S^*
}.
\tag{241}
\end{equation}

The relation

\begin{equation}
A^H=S^T
\tag{242}
\end{equation}

follows from

\begin{equation}
(S^*)^H=S^T.
\tag{243}
\end{equation}

This mapping permits the standard spectral initialization procedure to be applied without changing the definition of the channel vector $\mathbf g_i$.

\subsection{Augmented Channel Vector}

The received signal contains one unknown component, \(A^H\mathbf g_i\), and one known component, \(\mathbf b_i\):

\begin{equation}
A^H\mathbf g_i+\mathbf b_i.
\tag{244}
\end{equation}

The unknown channel and the known reference are combined into one augmented vector:

\begin{equation}
\boxed{
\bar{\mathbf g}_i
=
\begin{bmatrix}
\mathbf g_i\\
1
\end{bmatrix}
\in
\mathbb C^{(K+1)\times1}
}.
\tag{245}
\end{equation}

The last element is fixed to one so that

\begin{equation}
\begin{bmatrix}
A^H & \mathbf b_i
\end{bmatrix}
\begin{bmatrix}
\mathbf g_i\\
1
\end{bmatrix}
=
A^H\mathbf g_i+\mathbf b_i.
\tag{246}
\end{equation}

Define the augmented observation matrix as

\begin{equation}
\boxed{
\bar A_i
=
\begin{bmatrix}
A^H & \mathbf b_i
\end{bmatrix}
}.
\tag{247}
\end{equation}

Since

\begin{equation}
A^H=S^T,
\tag{248}
\end{equation}

we have

\begin{equation}
\boxed{
\bar A_i
=
\begin{bmatrix}
S^T & \mathbf b_i
\end{bmatrix}
}.
\tag{249}
\end{equation}

The dimensions are

\begin{equation}
S^T\in\mathbb C^{P\times K},
\qquad
\mathbf b_i\in\mathbb C^{P\times1},
\tag{250}
\end{equation}

and hence

\begin{equation}
\boxed{
\bar A_i\in\mathbb C^{P\times(K+1)}
}.
\tag{251}
\end{equation}

The measurement model becomes

\begin{equation}
\boxed{
\mathbf y_i
=
\left|
\bar A_i\bar{\mathbf g}_i+\mathbf w_i
\right|
}.
\tag{252}
\end{equation}

\subsection{Construction of the Augmented Observation Vectors}

The Hermitian transpose of the augmented matrix is

\begin{equation}
\bar A_i^H
\in
\mathbb C^{(K+1)\times P}.
\tag{253}
\end{equation}

Consider its $p$-th column:

\begin{equation}
\boxed{
\bar{\mathbf a}_{i,p}
=
\begin{bmatrix}
s_{1,p}^{*}\\
s_{2,p}^{*}\\
\vdots\\
s_{K,p}^{*}\\
b_{i,p}^{*}
\end{bmatrix}
\in
\mathbb C^{(K+1)\times1}
}.
\tag{254}
\end{equation}

Since

\begin{equation}
\bar A_i^H
=
\begin{bmatrix}
\bar{\mathbf a}_{i,1}
&
\bar{\mathbf a}_{i,2}
&
\cdots
&
\bar{\mathbf a}_{i,P}
\end{bmatrix},
\tag{255}
\end{equation}

the corresponding inner product is

\begin{equation}
\bar{\mathbf a}_{i,p}^{H}
\bar{\mathbf g}_i
=
\sum_{k=1}^{K}
s_{k,p}g_{i,k}
+
b_{i,p}.
\tag{256}
\end{equation}

Thus,

\begin{equation}
\boxed{
\bar{\mathbf a}_{i,p}^{H}
\bar{\mathbf g}_i
=
\left[S^T\mathbf g_i\right]_p+b_{i,p}
}.
\tag{257}
\end{equation}

\section{Spectral Matrix Construction}

\subsection{Weighted Covariance Matrix}

The spectral method constructs the weighted covariance matrix

\begin{equation}
\boxed{
M_i
=
\sum_{p=1}^{P}
y_{i,p}
\bar{\mathbf a}_{i,p}
\bar{\mathbf a}_{i,p}^{H}
}.
\tag{258}
\end{equation}

The dimensions are

\begin{equation}
\bar{\mathbf a}_{i,p}
\in
\mathbb C^{(K+1)\times1},
\tag{259}
\end{equation}

and therefore

\begin{equation}
\bar{\mathbf a}_{i,p}
\bar{\mathbf a}_{i,p}^{H}
\in
\mathbb C^{(K+1)\times(K+1)}.
\tag{260}
\end{equation}

Consequently,

\begin{equation}
\boxed{
M_i
\in
\mathbb C^{(K+1)\times(K+1)}
}.
\tag{261}
\end{equation}

The spectral matrix can also be written as

\begin{equation}
M_i
=
\sum_{p=1}^{P}
y_{i,p}
\bar{\mathbf a}_{i,p}
\bar{\mathbf a}_{i,p}^{H}.
\tag{262}
\end{equation}

This matrix contains the magnitude measurements $y_{i,p}$ as weights for the corresponding observation-vector outer products.

\subsection{Principal Eigenvector}

Perform the eigendecomposition

\begin{equation}
\boxed{
M_i\mathbf v_{i,r}
=
\lambda_{i,r}\mathbf v_{i,r}
}.
\tag{263}
\end{equation}

Select the eigenvector associated with the largest eigenvalue:

\begin{equation}
\boxed{
\mathbf v_i
=
\mathbf v_{i,r_{\max}}
}.
\tag{264}
\end{equation}

Equivalently,

\begin{equation}
M_i\mathbf v_i
=
\lambda_{\max}(M_i)\mathbf v_i.
\tag{265}
\end{equation}

The principal eigenvector satisfies

\begin{equation}
\boxed{
\mathbf v_i
\in
\mathbb C^{(K+1)\times1}
}.
\tag{266}
\end{equation}

The spectral method uses the principal eigenvector as an estimate of the direction of the augmented channel vector.

Thus,

\begin{equation}
\mathbf v_i
\approx
e^{j\phi_i}
\frac{\bar{\mathbf g}_i}
{\|\bar{\mathbf g}_i\|_2},
\tag{267}
\end{equation}

where $\phi_i$ represents an arbitrary phase.

The eigenvector provides the direction of the augmented channel, but its magnitude is not yet determined.

\section{Spectral Magnitude Estimation}

\subsection{Magnitude Estimation}

Assume the initial augmented channel vector has the form

\begin{equation}
\boxed{
\bar{\mathbf g}_i^{(0)}
=
r_i\mathbf v_i
},
\tag{268}
\end{equation}

where $r_i$ is a real nonnegative scalar.

The predicted magnitude corresponding to this estimate is

\begin{equation}
\left|
\bar A_i\bar{\mathbf g}_i^{(0)}
\right|
=
\left|
r_i\bar A_i\mathbf v_i
\right|.
\tag{269}
\end{equation}

Since $r_i\geq0$,

\begin{equation}
\left|
r_i\bar A_i\mathbf v_i
\right|
=
r_i
\left|
\bar A_i\mathbf v_i
\right|.
\tag{270}
\end{equation}

The magnitude should approximate the measured vector $\mathbf y_i$.

Therefore,

\begin{equation}
\boxed{
r_i
=
\arg\min_{r_i}
\left\|
r_i
\left|
\bar A_i\mathbf v_i
\right|
-
\mathbf y_i
\right\|_2^2
}.
\tag{271}
\end{equation}

Define

\begin{equation}
\mathbf q_i
=
\left|
\bar A_i\mathbf v_i
\right|.
\tag{272}
\end{equation}

Then the scalar least-squares problem becomes

\begin{equation}
r_i
=
\arg\min_{r_i}
\left\|
r_i\mathbf q_i-\mathbf y_i
\right\|_2^2.
\tag{273}
\end{equation}

Expanding the squared norm,

\begin{equation}
\left\|
r_i\mathbf q_i-\mathbf y_i
\right\|_2^2
=
r_i^2\mathbf q_i^T\mathbf q_i
-
2r_i\mathbf q_i^T\mathbf y_i
+
\mathbf y_i^T\mathbf y_i.
\tag{274}
\end{equation}

Since $r_i$ is real,

\begin{equation}
\frac{\partial}{\partial r_i}
\left\|
r_i\mathbf q_i-\mathbf y_i
\right\|_2^2
=
2r_i\|\mathbf q_i\|_2^2
-
2\mathbf q_i^T\mathbf y_i.
\tag{275}
\end{equation}

Setting the derivative equal to zero,

\begin{equation}
2r_i\|\mathbf q_i\|_2^2
-
2\mathbf q_i^T\mathbf y_i
=
0.
\tag{276}
\end{equation}

Therefore,

\begin{equation}
r_i\|\mathbf q_i\|_2^2
=
\mathbf q_i^T\mathbf y_i,
\tag{277}
\end{equation}

and

\begin{equation}
\boxed{
r_i
=
\frac{
\mathbf q_i^T\mathbf y_i
}{
\|\mathbf q_i\|_2^2
}
}.
\tag{278}
\end{equation}

Using (272),

\begin{equation}
\boxed{
r_i
=
\frac{
\left|
\bar A_i\mathbf v_i
\right|^T
\mathbf y_i
}{
\left\|
\bar A_i\mathbf v_i
\right\|_2^2
}
}.
\tag{279}
\end{equation}

The initial augmented channel estimate is therefore

\begin{equation}
\boxed{
\bar{\mathbf g}_i^{(0)}
=
r_i\mathbf v_i
}.
\tag{280}
\end{equation}

\section{Global Phase Removal}

\subsection{Arbitrary Global Phase}

Magnitude measurements cannot distinguish between

\begin{equation}
\bar{\mathbf g}_i
\tag{281}
\end{equation}

and

\begin{equation}
e^{j\phi_i}\bar{\mathbf g}_i,
\tag{282}
\end{equation}

because

\begin{equation}
\left|
\bar A_i
e^{j\phi_i}
\bar{\mathbf g}_i
\right|
=
\left|
e^{j\phi_i}
\bar A_i
\bar{\mathbf g}_i
\right|
=
\left|
\bar A_i\bar{\mathbf g}_i
\right|.
\tag{283}
\end{equation}

However, the augmented vector is defined as

\begin{equation}
\bar{\mathbf g}_i
=
\begin{bmatrix}
\mathbf g_i\\
1
\end{bmatrix},
\tag{284}
\end{equation}

so the final element must have zero phase.

Let

\begin{equation}
\boxed{
\phi_i
=
\angle
\left(
[\bar{\mathbf g}_i^{(0)}]_{K+1}
\right)
}.
\tag{285}
\end{equation}

Multiply the entire augmented vector by

\begin{equation}
e^{-j\phi_i}.
\tag{286}
\end{equation}

Thus,

\begin{equation}
\boxed{
\widetilde{\mathbf g}_i^{(0)}
=
e^{-j\phi_i}
\bar{\mathbf g}_i^{(0)}
}.
\tag{287}
\end{equation}

The last element becomes

\begin{equation}
[\widetilde{\mathbf g}_i^{(0)}]_{K+1}
=
e^{-j\phi_i}
[\bar{\mathbf g}_i^{(0)}]_{K+1}.
\tag{288}
\end{equation}

Using (285),

\begin{equation}
[\bar{\mathbf g}_i^{(0)}]_{K+1}
=
\left|
[\bar{\mathbf g}_i^{(0)}]_{K+1}
\right|
e^{j\phi_i},
\tag{289}
\end{equation}

and therefore,

\begin{align}
[\widetilde{\mathbf g}_i^{(0)}]_{K+1}
&=
e^{-j\phi_i}
\left|
[\bar{\mathbf g}_i^{(0)}]_{K+1}
\right|
e^{j\phi_i}
\nonumber\\
&=
\left|
[\bar{\mathbf g}_i^{(0)}]_{K+1}
\right|.
\tag{290}
\end{align}

Hence, the final element has zero phase.

The channel estimate for the $i$-th antenna is obtained by retaining the first $K$ elements:

\begin{equation}
\boxed{
\mathbf g_i^{(0)}
=
\widetilde{\mathbf g}_i^{(0)}(1:K)
}.
\tag{291}
\end{equation}

\subsection{Construction of the Initial Channel Matrix}

The above spectral initialization is performed independently for all $I$ vapor cells.

For

\begin{equation}
i=1,2,\ldots,I,
\tag{292}
\end{equation}

obtain

\begin{equation}
\mathbf g_i^{(0)}
\in
\mathbb C^{K\times1}.
\tag{293}
\end{equation}

The initial channel matrix is constructed by placing the channel vectors as rows:

\begin{equation}
\boxed{
G^{(0)}
=
\begin{bmatrix}
(\mathbf g_1^{(0)})^T\\
(\mathbf g_2^{(0)})^T\\
\vdots\\
(\mathbf g_I^{(0)})^T
\end{bmatrix}
\in
\mathbb C^{I\times K}
}.
\tag{294}
\end{equation}

This matrix provides the initial channel estimate for the subsequent GS iterations.

\section{Gerchberg--Saxton Iteration}

\subsection{Step 1: Prediction of the Complex Received Signal}

Suppose that the channel estimate from the previous iteration is

\begin{equation}
G^{(t-1)}
\in
\mathbb C^{I\times K}.
\tag{295}
\end{equation}

For the $i$-th antenna,

\begin{equation}
\mathbf g_i^{(t-1)}
\in
\mathbb C^{K\times1}.
\tag{296}
\end{equation}

The predicted complex received signal is

\begin{equation}
\boxed{
\mathbf x_i^{(t)}
=
S^T\mathbf g_i^{(t-1)}
+
\mathbf b_i
}.
\tag{297}
\end{equation}

The noise realization is not added to the predicted signal because the actual noise realization is unknown. The measured magnitude already contains the noise contribution.

The dimensions are

\begin{equation}
S^T\in\mathbb C^{P\times K},
\qquad
\mathbf g_i^{(t-1)}
\in
\mathbb C^{K\times1},
\tag{298}
\end{equation}

hence

\begin{equation}
\mathbf x_i^{(t)}
\in
\mathbb C^{P\times1}.
\tag{299}
\end{equation}

\subsection{Step 2: Phase Estimation}

The phase of the predicted complex received signal is

\begin{equation}
\boxed{
\boldsymbol{\theta}_i^{(t)}
=
\angle
\left(
S^T\mathbf g_i^{(t-1)}
+
\mathbf b_i
\right)
}.
\tag{300}
\end{equation}

Collecting the phase values gives

\begin{equation}
\boldsymbol{\theta}_i^{(t)}
=
\begin{bmatrix}
\theta_{i,1}^{(t)}\\
\theta_{i,2}^{(t)}\\
\vdots\\
\theta_{i,P}^{(t)}
\end{bmatrix}
\in
\mathbb R^{P\times1}.
\tag{301}
\end{equation}

The corresponding unit-magnitude phase vector is

\begin{equation}
e^{j\boldsymbol{\theta}_i^{(t)}}.
\tag{302}
\end{equation}

Each element has unit magnitude:

\begin{equation}
\left|
e^{j\theta_{i,p}^{(t)}}
\right|
=
1.
\tag{303}
\end{equation}

\subsection{Step 3: Replace the Predicted Magnitude with the Measured Magnitude}

The measured vector is

\begin{equation}
\mathbf y_i
=
|\mathbf x_i|.
\tag{304}
\end{equation}

The GS procedure retains the measured magnitude and replaces only the unknown phase by the current phase estimate.

Define

\begin{equation}
\boxed{
\mathbf r_i^{(t)}
=
\mathbf y_i
\circ
e^{j\boldsymbol{\theta}_i^{(t)}}
}.
\tag{305}
\end{equation}

Element-wise,

\begin{equation}
\boxed{
r_{i,p}^{(t)}
=
y_{i,p}
e^{j\theta_{i,p}^{(t)}}
}.
\tag{306}
\end{equation}

Therefore,

\begin{equation}
\left|
r_{i,p}^{(t)}
\right|
=
y_{i,p}
\left|
e^{j\theta_{i,p}^{(t)}}
\right|
=
y_{i,p}.
\tag{307}
\end{equation}

Thus,

\begin{equation}
\boxed{
|\mathbf r_i^{(t)}|
=
\mathbf y_i
}.
\tag{308}
\end{equation}

This operation replaces the unknown complex received signal with a complex vector having the measured magnitude and the currently estimated phase.

\section{Reference Removal and Least-Squares Channel Update}

\subsection{Step 4: Remove the Known Reference}

The model for the reconstructed received signal is

\begin{equation}
\mathbf r_i^{(t)}
\approx
S^T\mathbf g_i+\mathbf b_i.
\tag{309}
\end{equation}

Therefore,

\begin{equation}
\boxed{
\mathbf r_i^{(t)}-\mathbf b_i
\approx
S^T\mathbf g_i
}.
\tag{310}
\end{equation}

After removing the known reference, the channel estimation becomes a standard linear least-squares problem.

\subsection{Step 5: Least-Squares Channel Estimation}

The channel vector is obtained from

\begin{equation}
\boxed{
\mathbf g_i^{(t)}
=
\arg\min_{\mathbf g}
\left\|
\mathbf r_i^{(t)}
-
\mathbf b_i
-
S^T\mathbf g
\right\|_2^2
}.
\tag{311}
\end{equation}

To obtain the standard least-squares form, define

\begin{equation}
A=S^*.
\tag{312}
\end{equation}

Then

\begin{equation}
A^H=S^T.
\tag{313}
\end{equation}

Therefore,

\begin{equation}
\mathbf g_i^{(t)}
=
\arg\min_{\mathbf g}
\left\|
(\mathbf r_i^{(t)}-\mathbf b_i)
-
A^H\mathbf g
\right\|_2^2.
\tag{314}
\end{equation}

The standard least-squares solution is

\begin{equation}
\boxed{
\mathbf g_i^{(t)}
=
(AA^H)^{-1}
A
\left(
\mathbf r_i^{(t)}-\mathbf b_i
\right)
}.
\tag{315}
\end{equation}

This solution requires \(AA^H=S^*S^T\in\mathbb C^{K\times K}\) to be invertible, i.e.
\(S\) must have full row rank \(K\), which needs \(P\geq K\).

Since

\begin{equation}
A=S^*,
\tag{316}
\end{equation}

and

\begin{equation}
A^H=S^T,
\tag{317}
\end{equation}

we have

\begin{equation}
AA^H
=
S^*S^T.
\tag{318}
\end{equation}

Hence,

\begin{equation}
\boxed{
\mathbf g_i^{(t)}
=
(S^*S^T)^{-1}
S^*
\left(
\mathbf r_i^{(t)}-\mathbf b_i
\right)
}.
\tag{319}
\end{equation}

Using (305),

\begin{equation}
\boxed{
\mathbf g_i^{(t)}
=
(S^*S^T)^{-1}
S^*
\left(
\mathbf y_i
\circ
e^{j\boldsymbol{\theta}_i^{(t)}}
-
\mathbf b_i
\right)
}.
\tag{320}
\end{equation}

\section{Matrix Form of the GS Iteration}

\subsection{Simultaneous Update for All Antennas}

Instead of processing each antenna separately, define

\begin{equation}
\boxed{
G^{(t-1)}
\in
\mathbb C^{I\times K}
}.
\tag{321}
\end{equation}

The predicted complex received matrix is

\begin{equation}
\boxed{
X^{(t)}
=
G^{(t-1)}S+B
}.
\tag{322}
\end{equation}

The dimensions are

\begin{equation}
G^{(t-1)}S:
(I\times K)(K\times P)
=
I\times P,
\tag{323}
\end{equation}

and

\begin{equation}
B\in\mathbb C^{I\times P}.
\tag{324}
\end{equation}

Hence,

\begin{equation}
\boxed{
X^{(t)}
\in
\mathbb C^{I\times P}
}.
\tag{325}
\end{equation}

The phase matrix is

\begin{equation}
\boxed{
\Theta^{(t)}
=
\angle X^{(t)}
}.
\tag{326}
\end{equation}

The reconstructed complex measurement matrix is

\begin{equation}
\boxed{
R^{(t)}
=
Y\circ
e^{j\Theta^{(t)}}
}.
\tag{327}
\end{equation}

Since $Y$ contains the measured magnitudes,

\begin{equation}
|R^{(t)}|
=
Y.
\tag{328}
\end{equation}

The known reference matrix is then removed:

\begin{equation}
\boxed{
R^{(t)}-B
}.
\tag{329}
\end{equation}

The channel update is obtained by solving

\begin{equation}
\boxed{
G^{(t)}
=
\arg\min_G
\left\|
R^{(t)}-B-GS
\right\|_F^2
}.
\tag{330}
\end{equation}

The least-squares solution is

\begin{equation}
\boxed{
G^{(t)}
=
(R^{(t)}-B)
S^H
(SS^H)^{-1}
}.
\tag{331}
\end{equation}

The dimensions are

\begin{equation}
R^{(t)}-B
\in
\mathbb C^{I\times P},
\tag{332}
\end{equation}

\begin{equation}
S^H
\in
\mathbb C^{P\times K},
\tag{333}
\end{equation}

and

\begin{equation}
SS^H
\in
\mathbb C^{K\times K}.
\tag{334}
\end{equation}

Therefore,

\begin{equation}
(R^{(t)}-B)S^H
\in
\mathbb C^{I\times K},
\tag{335}
\end{equation}

and

\begin{equation}
\boxed{
G^{(t)}
\in
\mathbb C^{I\times K}
}.
\tag{336}
\end{equation}

\subsection{Equivalence with the Per-Antenna LS Solution}

For the $i$-th row of the matrix model,

\begin{equation}
[R^{(t)}-B]_{i,:}
=
(\mathbf r_i^{(t)}-\mathbf b_i)^T.
\tag{337}
\end{equation}

The corresponding row of $G^{(t)}$ is

\begin{equation}
[G^{(t)}]_{i,:}
=
(\mathbf g_i^{(t)})^T.
\tag{338}
\end{equation}

From (331),

\begin{equation}
[G^{(t)}]_{i,:}
=
[R^{(t)}-B]_{i,:}
S^H
(SS^H)^{-1}.
\tag{339}
\end{equation}

Taking the transpose gives the corresponding channel-vector representation,

\begin{equation}
\mathbf g_i^{(t)}
=
(S^*S^T)^{-1}
S^*
(\mathbf r_i^{(t)}-\mathbf b_i).
\tag{340}
\end{equation}

Thus, the matrix update and the per-antenna least-squares solution represent the same channel-estimation operation under different orientations of the channel vectors.

\section{Convergence Criterion}

\subsection{Difference Between Consecutive Channel Estimates}

After obtaining $G^{(t)}$, compare it with the channel estimate from the previous iteration:

\begin{equation}
\boxed{
\Delta_t
=
\left\|
G^{(t)}
-
G^{(t-1)}
\right\|_F^2
}.
\tag{341}
\end{equation}

If

\begin{equation}
\boxed{
\Delta_t<\epsilon
},
\tag{342}
\end{equation}

where $\epsilon$ is a prescribed convergence tolerance, the iteration is terminated.

Otherwise, the GS iterations continue.

A maximum iteration number may also be imposed as a stopping condition:

\begin{equation}
\boxed{
t\geq T_{\max}
}.
\tag{343}
\end{equation}

Therefore, the algorithm terminates when either

\begin{equation}
\Delta_t<\epsilon
\tag{344}
\end{equation}

or

\begin{equation}
t\geq T_{\max}.
\tag{345}
\end{equation}

\section{Final Biased GS Algorithm for the 1D Array}

The complete biased GS channel-estimation algorithm consists of spectral initialization followed by iterative phase estimation and least-squares channel updates.

\subsection{Algorithm Input}

The required inputs are

\begin{equation}
\boxed{
Y,\quad S,\quad B,\quad T_{\max},\quad \epsilon
}.
\tag{346}
\end{equation}

For the per-antenna spectral initialization,

\begin{equation}
\mathbf y_i
\in
\mathbb R^{P\times1},
\tag{347}
\end{equation}

\begin{equation}
S
\in
\mathbb C^{K\times P},
\tag{348}
\end{equation}

and

\begin{equation}
\mathbf b_i
\in
\mathbb C^{P\times1}.
\tag{349}
\end{equation}

\subsection{Spectral Initialization Procedure}

For

\begin{equation}
i=1,2,\ldots,I,
\tag{350}
\end{equation}

construct

\begin{equation}
\boxed{
\bar A_i
=
\begin{bmatrix}
S^T & \mathbf b_i
\end{bmatrix}
}.
\tag{351}
\end{equation}

Let $\bar{\mathbf a}_{i,p}$ denote the $p$-th column of $\bar A_i^H$.

Construct

\begin{equation}
\boxed{
M_i
=
\sum_{p=1}^{P}
y_{i,p}
\bar{\mathbf a}_{i,p}
\bar{\mathbf a}_{i,p}^{H}
}.
\tag{352}
\end{equation}

Obtain the principal eigenvector

\begin{equation}
\boxed{
\mathbf v_i
=
\operatorname{PrincipalEigenvector}(M_i)
}.
\tag{353}
\end{equation}

Estimate the magnitude scaling factor:

\begin{equation}
\boxed{
r_i
=
\frac{
|\bar A_i\mathbf v_i|^T
\mathbf y_i
}{
\|\bar A_i\mathbf v_i\|_2^2
}
}.
\tag{354}
\end{equation}

Construct the initial augmented channel vector:

\begin{equation}
\boxed{
\bar{\mathbf g}_i^{(0)}
=
r_i\mathbf v_i
}.
\tag{355}
\end{equation}

Remove its arbitrary global phase:

\begin{equation}
\boxed{
\widetilde{\mathbf g}_i^{(0)}
=
e^{-j
\angle
([\bar{\mathbf g}_i^{(0)}]_{K+1})}
\bar{\mathbf g}_i^{(0)}
}.
\tag{356}
\end{equation}

Retain the first $K$ elements:

\begin{equation}
\boxed{
\mathbf g_i^{(0)}
=
\widetilde{\mathbf g}_i^{(0)}(1:K)
}.
\tag{357}
\end{equation}

After processing all antennas, construct

\begin{equation}
\boxed{
G^{(0)}
=
\begin{bmatrix}
(\mathbf g_1^{(0)})^T\\
(\mathbf g_2^{(0)})^T\\
\vdots\\
(\mathbf g_I^{(0)})^T
\end{bmatrix}
}.
\tag{358}
\end{equation}

\subsection{GS Iterations}

For

\begin{equation}
t=1,2,\ldots,T_{\max},
\tag{359}
\end{equation}

calculate the predicted complex received matrix:

\begin{equation}
\boxed{
X^{(t)}
=
G^{(t-1)}S+B
}.
\tag{360}
\end{equation}

Calculate its phase:

\begin{equation}
\boxed{
\Theta^{(t)}
=
\angle X^{(t)}
}.
\tag{361}
\end{equation}

Replace the predicted magnitude by the measured magnitude:

\begin{equation}
\boxed{
R^{(t)}
=
Y\circ e^{j\Theta^{(t)}}
}.
\tag{362}
\end{equation}

Remove the reference:

\begin{equation}
\boxed{
R^{(t)}-B
}.
\tag{363}
\end{equation}

Update the channel using least squares:

\begin{equation}
\boxed{
G^{(t)}
=
(R^{(t)}-B)
S^H
(SS^H)^{-1}
}.
\tag{364}
\end{equation}

Evaluate the convergence measure:

\begin{equation}
\boxed{
\Delta_t
=
\left\|
G^{(t)}
-
G^{(t-1)}
\right\|_F^2
}.
\tag{365}
\end{equation}

If

\begin{equation}
\boxed{
\Delta_t<\epsilon,
}
\tag{366}
\end{equation}

terminate the iterations.

Otherwise, continue until

\begin{equation}
t=T_{\max}.
\tag{367}
\end{equation}

The final channel estimate is

\begin{equation}
\boxed{
\widehat G=G^{(t)}
}.
\tag{368}
\end{equation}

\clearpage
\section{Complete Algorithm}

The complete biased GS algorithm for the 1D atomic receiver array can therefore be summarized as

\begin{equation}
\boxed{
\begin{array}{l}
\textbf{Input: }
Y,\;S,\;B,\;T_{\max},\;\epsilon
\\[2mm]
\textbf{Spectral Initialization:}
\\
\qquad
\text{For }i=1,2,\ldots,I:
\\
\qquad
\bar A_i=[S^T\;\;\mathbf b_i]
\\
\qquad
\bar{\mathbf a}_{i,p}
=
\text{$p$th column of }\bar A_i^H
\\
\qquad
M_i
=
\displaystyle\sum_{p=1}^{P}
y_{i,p}
\bar{\mathbf a}_{i,p}
\bar{\mathbf a}_{i,p}^{H}
\\
\qquad
\mathbf v_i
=
\text{principal eigenvector of }M_i
\\
\qquad
r_i
=
\displaystyle
\frac{
|\bar A_i\mathbf v_i|^T\mathbf y_i
}{
\|\bar A_i\mathbf v_i\|_2^2
}
\\
\qquad
\bar{\mathbf g}_i^{(0)}
=
r_i\mathbf v_i
\\
\qquad
\widetilde{\mathbf g}_i^{(0)}
=
e^{-j\angle([\bar{\mathbf g}_i^{(0)}]_{K+1})}
\bar{\mathbf g}_i^{(0)}
\\
\qquad
\mathbf g_i^{(0)}
=
\widetilde{\mathbf g}_i^{(0)}(1:K)
\\
\qquad
\text{End}
\\[2mm]
\qquad
G^{(0)}
=
\begin{bmatrix}
(\mathbf g_1^{(0)})^T\\
\vdots\\
(\mathbf g_I^{(0)})^T
\end{bmatrix}
\\[3mm]
\textbf{GS Iterations:}
\\
\qquad
\text{For }t=1,2,\ldots,T_{\max}:
\\
\qquad
X^{(t)}
=
G^{(t-1)}S+B
\\
\qquad
\Theta^{(t)}
=
\angle X^{(t)}
\\
\qquad
R^{(t)}
=
Y\circ e^{j\Theta^{(t)}}
\\
\qquad
G^{(t)}
=
(R^{(t)}-B)
S^H
(SS^H)^{-1}
\\
\qquad
\Delta_t
=
\|G^{(t)}-G^{(t-1)}\|_F^2
\\
\qquad
\text{If }\Delta_t<\epsilon:
\text{ stop}
\\
\qquad
\text{End}
\\[2mm]
\textbf{Output: }
\boxed{\widehat G=G^{(t)}}.
\end{array}
}
\tag{369}
\end{equation}

% ============================================================
% END OF PAGES 41--51
% ============================================================

% ============================================================
% 2D ANTENNA ARRAY, PGD, AND CRLB
% ============================================================

\section{Extension to a 2D Antenna Array}

The one-dimensional channel model is extended to a planar array of
vapor cells containing $I_1\times I_2$ receiver elements. The two
array indices are denoted by
\[
i_1\in\{1,\ldots,I_1\},
\qquad
i_2\in\{1,\ldots,I_2\}.
\]
The spatial phase of the $l$-th path of user $k$ is determined
independently along the two array dimensions, in the same way as the
1D phase (8) was obtained along one dimension.

\subsection{Two-Dimensional Spatial Phase}

Let $d_1$ and $d_2$ denote the inter-element spacings along the two
array axes. For the $l$-th path of user $k$, let $\theta_{l,k}$ and
$\phi_{l,k}$ denote the elevation and azimuth angles, respectively.
The phase increments between adjacent elements along the two
dimensions are
\begin{equation}
\boxed{
 u_{l,k}=\frac{2\pi d_1}{\lambda}\cos\theta_{l,k},
 \qquad
 v_{l,k}=\frac{2\pi d_2}{\lambda}
 \sin\theta_{l,k}\cos\phi_{l,k}.
}
\tag{370}
\end{equation}

Both $u_{l,k}$ and $v_{l,k}$ are real scalars. Taking element $(1,1)$
as the phase reference, the total spatial phase at the $(i_1,i_2)$-th
receiver element is
\begin{equation}
\boxed{
 (i_1-1)u_{l,k}+(i_2-1)v_{l,k}.
}
\tag{371}
\end{equation}

Using $e^{-j(a+b)}=e^{-ja}e^{-jb}$,
\begin{equation}
 e^{-j[(i_1-1)u_{l,k}+(i_2-1)v_{l,k}]}
 =e^{-j(i_1-1)u_{l,k}}
  e^{-j(i_2-1)v_{l,k}}.
\tag{372}
\end{equation}

The phase factor therefore separates into a term depending only on
$i_1$ and a term depending only on $i_2$. This separability gives the
2D channel matrix its low-rank structure.

\subsection{Two-Dimensional Channel Coefficient}

Extending (75) to two array dimensions, the path-dependent coupling
coefficient of the $l$-th path of user $k$ is defined as
\begin{equation}
\boxed{
 \beta_{l,k}
 =
 \frac{1}{\hbar}
 \boldsymbol{\mu}_{eg}^{T}
 \boldsymbol{\epsilon}_{k,l}
 \alpha_{l,k}
 \in\mathbb C.
}
\tag{373}
\end{equation}

For the rank-constrained model, the polarization of each path is
assumed to be invariant across the array aperture,
\begin{equation}
\boxed{
 \boldsymbol{\epsilon}_{i_1,i_2,k,l}
 =
 \boldsymbol{\epsilon}_{k,l}.
}
\tag{374}
\end{equation}

Under this assumption, the channel coefficient between user $k$ and the
$(i_1,i_2)$-th receiver element is
\begin{equation}
\boxed{
 g_{i_1,i_2,k}
 =
 \sum_{l=1}^{L_k}
 \beta_{l,k}
 e^{-j[(i_1-1)u_{l,k}+(i_2-1)v_{l,k}]}
 \in\mathbb C.
}
\tag{375}
\end{equation}

Define the two spatial steering vectors
\begin{equation}
\mathbf a_{1,l,k}
 =
 \begin{bmatrix}
 1 & e^{-ju_{l,k}} & \cdots & e^{-j(I_1-1)u_{l,k}}
 \end{bmatrix}^{T}
 \in\mathbb C^{I_1\times1},
\tag{376}
\end{equation}
\begin{equation}
\mathbf a_{2,l,k}
 =
 \begin{bmatrix}
 1 & e^{-jv_{l,k}} & \cdots & e^{-j(I_2-1)v_{l,k}}
 \end{bmatrix}^{T}
 \in\mathbb C^{I_2\times1}.
\tag{377}
\end{equation}

The $(i_1,i_2)$-th element of the outer product
$\mathbf a_{1,l,k}\mathbf a_{2,l,k}^{T}$ is
\[
[\mathbf a_{1,l,k}]_{i_1}[\mathbf a_{2,l,k}]_{i_2}
=
e^{-j(i_1-1)u_{l,k}}e^{-j(i_2-1)v_{l,k}},
\]
which is exactly the phase factor in (372). Hence the contribution of
the $l$-th path to the $k$-th user's channel matrix is
\begin{equation}
\mathbf G_k^{(l)}
 =
 \beta_{l,k}
 \mathbf a_{1,l,k}\mathbf a_{2,l,k}^{T},
 \qquad
 \underbrace{\mathbf a_{1,l,k}}_{I_1\times1}
 \underbrace{\mathbf a_{2,l,k}^{T}}_{1\times I_2}
 =
 \underbrace{\mathbf G_k^{(l)}}_{I_1\times I_2}.
\tag{378}
\end{equation}

Since both steering vectors are nonzero (their first entries equal
one) and $\beta_{l,k}\neq0$, every column of $\mathbf G_k^{(l)}$ is a
multiple of $\mathbf a_{1,l,k}$. Therefore,
\begin{equation}
\operatorname{rank}\!\left(
\mathbf G_k^{(l)}
\right)=1.
\tag{379}
\end{equation}

The complete channel matrix of user $k$ is the superposition of its
$L_k$ paths:
\begin{equation}
\boxed{
\mathbf G_k
 =
 \sum_{l=1}^{L_k}\mathbf G_k^{(l)}
 \in\mathbb C^{I_1\times I_2},
 \qquad
 [\mathbf G_k]_{i_1,i_2}=g_{i_1,i_2,k}.
}
\tag{380}
\end{equation}

Using the rank subadditivity property
$\operatorname{rank}(X+Y)\leq\operatorname{rank}(X)+\operatorname{rank}(Y)$,
\begin{equation}
\begin{aligned}
\operatorname{rank}(\mathbf G_k)
&=
\operatorname{rank}\!\left(
\sum_{l=1}^{L_k}\mathbf G_k^{(l)}
\right)
\leq
\sum_{l=1}^{L_k}
\operatorname{rank}\!\left(\mathbf G_k^{(l)}\right)
=
\sum_{l=1}^{L_k}1
=L_k.
\end{aligned}
\tag{381}
\end{equation}

Since the rank of an $I_1\times I_2$ matrix also cannot exceed
$\min(I_1,I_2)$,
\begin{equation}
\boxed{
\operatorname{rank}(\mathbf G_k)\leq \min(L_k,I_1,I_2).
}
\tag{382}
\end{equation}

The assumption in (374) is required for this rank argument. If the
polarization varies with $(i_1,i_2)$, the path matrix becomes
$\alpha_{l,k}\,\mathbf C_{k,l}\circ(\mathbf a_{1,l,k}\mathbf a_{2,l,k}^{T})$,
where $[\mathbf C_{k,l}]_{i_1,i_2}=\hbar^{-1}\boldsymbol\mu_{eg}^{T}
\boldsymbol\epsilon_{i_1,i_2,k,l}$, and the rank-one decomposition in
(378) no longer holds. In simulation, $\boldsymbol\epsilon_{k,l}$ must
therefore be drawn once per path, not once per vapor cell.

\subsection{Two-Dimensional Reference Signal}

Extending (89) to two dimensions, the received reference signal at the
$(i_1,i_2)$-th receiver element and pilot instant $p$ is
\begin{equation}
\boxed{
 b_{i_1,i_2,p}
 =
 s_{b,p}
 \frac{1}{\hbar}
 \boldsymbol{\mu}_{eg}^{T}\boldsymbol{\epsilon}_{b,i_1,i_2}
 \alpha_b
 e^{-j[(i_1-1)u_b+(i_2-1)v_b]}
 \in\mathbb C,
}
\tag{383}
\end{equation}
where $u_b$ and $v_b$ are the phase increments of the reference path,
defined as in (370). As in the 1D case, the reference (LO) field is
treated as a plane wave across the array.

Consistent with (125)--(127), the reference is written in
magnitude--phase form,
\begin{equation}
 b_{i_1,i_2,p}=|b_{i_1,i_2,p}|e^{jz_{i_1,i_2,p}},
 \qquad
 z_{i_1,i_2,p}=\angle b_{i_1,i_2,p},
 \qquad
 \frac{|b_{i_1,i_2,p}|}{b_{i_1,i_2,p}}=e^{-jz_{i_1,i_2,p}}.
\tag{384}
\end{equation}

\subsection{Received Signal and Strong-Reference Linearisation}

As in (93), for pilot $p$ the received magnitude at receiver
$(i_1,i_2)$ is
\begin{equation}
 y_{i_1,i_2,p}
 =
 \left|
 \sum_{k=1}^{K}g_{i_1,i_2,k}s_{k,p}
 +b_{i_1,i_2,p}
 +n_{i_1,i_2,p}
 \right|
 \in\mathbb R.
\tag{385}
\end{equation}

Define the desired signal and the normalized perturbation
\begin{equation}
 a_{i_1,i_2,p}
 =
 \sum_{k=1}^{K}g_{i_1,i_2,k}s_{k,p},
 \qquad
 r_{i_1,i_2,p}
 =
 \frac{a_{i_1,i_2,p}+n_{i_1,i_2,p}}
 {b_{i_1,i_2,p}}.
\tag{386}
\end{equation}

Factoring $b_{i_1,i_2,p}$ out of (385) and using
$|1+r|^2=(1+r)(1+r^*)=1+2\operatorname{Re}(r)+|r|^2$, as in
(97)--(110),
\begin{equation}
\begin{aligned}
y_{i_1,i_2,p}
&=
|b_{i_1,i_2,p}|\,|1+r_{i_1,i_2,p}|\\
&=
|b_{i_1,i_2,p}|
\sqrt{
1+2\operatorname{Re}(r_{i_1,i_2,p})
+|r_{i_1,i_2,p}|^2
}.
\end{aligned}
\tag{387}
\end{equation}

Under the strong-reference condition $|r_{i_1,i_2,p}|\ll1$, the
second-order term $|r_{i_1,i_2,p}|^2$ is neglected and
$\sqrt{1+x}\approx1+x/2$ is used. Applying (384),
\begin{equation}
\begin{aligned}
y_{i_1,i_2,p}
&\approx
|b_{i_1,i_2,p}|
+|b_{i_1,i_2,p}|\operatorname{Re}(r_{i_1,i_2,p})\\
&=
|b_{i_1,i_2,p}|
+\operatorname{Re}\!\left[
 e^{-jz_{i_1,i_2,p}}
 \left(a_{i_1,i_2,p}+n_{i_1,i_2,p}\right)
\right],
\end{aligned}
\tag{388}
\end{equation}
where the real scalar $|b_{i_1,i_2,p}|$ was moved inside the real part.

Define the effective real-valued noise
\begin{equation}
\boxed{
\bar n_{i_1,i_2,p}
=
\operatorname{Re}\!\left(
 e^{-jz_{i_1,i_2,p}}
 n_{i_1,i_2,p}
\right)
\in\mathbb R.
}
\tag{389}
\end{equation}

If the complex noise is
\begin{equation}
n_{i_1,i_2,p}\sim\mathcal{CN}(0,\sigma^2),
\tag{390}
\end{equation}
then $e^{-jz_{i_1,i_2,p}}n_{i_1,i_2,p}\sim\mathcal{CN}(0,\sigma^2)$ by
circular symmetry, and its real part carries half of the variance:
\begin{equation}
\boxed{
\bar n_{i_1,i_2,p}
\sim
\mathcal N\!\left(0,\frac{\sigma^2}{2}\right).
}
\tag{391}
\end{equation}

The resulting linearised 2D observation model is
\begin{equation}
\boxed{
 y_{i_1,i_2,p}
 \approx
 \operatorname{Re}\!\left(
 e^{-jz_{i_1,i_2,p}}
 \sum_{k=1}^{K}g_{i_1,i_2,k}s_{k,p}
 \right)
 +|b_{i_1,i_2,p}|
 +\bar n_{i_1,i_2,p},
}
\tag{392}
\end{equation}
which is Eq.~(18) of Xu \emph{et al.} and the 2D counterpart of (134).

% ============================================================
\section{Tensor Representation and Mode-3 Unfolding}
% ============================================================

Define the channel tensor by stacking the $K$ user matrices
$\mathbf G_k$ along the third dimension:
\begin{equation}
\mathcal G\in\mathbb C^{I_1\times I_2\times K},
\qquad
[\mathcal G]_{i_1,i_2,k}=g_{i_1,i_2,k},
\qquad
\mathcal G(:,:,k)=\mathbf G_k.
\tag{393}
\end{equation}

Each frontal slice $\mathcal G(:,:,k)$ is therefore a channel matrix of
rank at most $L_k$ by (382). Similarly, define
\begin{equation}
\begin{aligned}
[\mathcal Y]_{i_1,i_2,p}&=y_{i_1,i_2,p},
&\mathcal Y&\in\mathbb R^{I_1\times I_2\times P},\\
[\mathcal Z]_{i_1,i_2,p}&=e^{-jz_{i_1,i_2,p}},
&\mathcal Z&\in\mathbb C^{I_1\times I_2\times P},\\
[\mathcal B]_{i_1,i_2,p}&=b_{i_1,i_2,p},
&\mathcal B&\in\mathbb C^{I_1\times I_2\times P},\\
[\mathcal N]_{i_1,i_2,p}&=\bar n_{i_1,i_2,p},
&\mathcal N&\in\mathbb R^{I_1\times I_2\times P}.
\end{aligned}
\tag{394}
\end{equation}

The pilot matrix is the same as in (145):
\begin{equation}
S\in\mathbb C^{K\times P},
\qquad
[S]_{k,p}=s_{k,p}.
\tag{395}
\end{equation}

The mode-3 product multiplies every mode-3 fibre
$\mathcal G(i_1,i_2,:)\in\mathbb C^{K}$ by $S^T\in\mathbb C^{P\times K}$:
\begin{equation}
\mathcal G\times_3 S^T
\in
\mathbb C^{I_1\times I_2\times P},
\qquad
[\mathcal G\times_3 S^T]_{i_1,i_2,p}
=\sum_{k=1}^{K}g_{i_1,i_2,k}s_{k,p}
=a_{i_1,i_2,p}.
\tag{396}
\end{equation}

Collecting (392) over all $(i_1,i_2,p)$, the linearised tensor model is
\begin{equation}
\boxed{
\mathcal Y
=
\operatorname{Re}\!\left[
\mathcal Z\circ
(\mathcal G\times_3 S^T)
\right]
+|\mathcal B|
+\mathcal N,
}
\tag{397}
\end{equation}
which is Eq.~(19) of Xu \emph{et al.}; here $\circ$ is the elementwise
product and $|\mathcal B|$ the elementwise magnitude.

Let the mode-3 unfolding place the mode-3 fibres as columns of a matrix
whose rows are indexed by the third index, with $(i_1,i_2)$ mapped to
the column index
\begin{equation}
q=i_1+(i_2-1)I_1,
\qquad
q=1,\ldots,I_1I_2,
\tag{398}
\end{equation}
which is column-major order. Then
\begin{equation}
G_{(3)}\in\mathbb C^{K\times I_1I_2},
\qquad
[G_{(3)}]_{k,q}=g_{i_1,i_2,k},
\tag{398a}
\end{equation}
and
\begin{equation}
Y_{(3)},N_{(3)}\in\mathbb R^{P\times I_1I_2},
\qquad
Z_{(3)},B_{(3)}\in\mathbb C^{P\times I_1I_2}.
\tag{399}
\end{equation}

For the mode-3 product, the unfolding identity
$(\mathcal X\times_3 U)_{(3)}=U\mathcal X_{(3)}$ gives
\begin{equation}
\boxed{
(\mathcal G\times_3 S^T)_{(3)}
=
S^T G_{(3)},
\qquad
\underbrace{S^T}_{P\times K}
\underbrace{G_{(3)}}_{K\times I_1I_2}
=
\underbrace{S^TG_{(3)}}_{P\times I_1I_2}.
}
\tag{400}
\end{equation}

Unfolding every term of (397) along mode 3,
\begin{equation}
\boxed{
Y_{(3)}
=
\operatorname{Re}\!\left[
Z_{(3)}\circ(S^T G_{(3)})
\right]
+|B_{(3)}|
+N_{(3)},
}
\tag{401}
\end{equation}
which is Eq.~(20) of Xu \emph{et al.} Compared with the 1D model
(163), the roles of rows and columns are transposed: rows now index
pilots and columns index antennas.

For user $k$, the $k$-th row of $G_{(3)}$ is
$\operatorname{vec}(\mathbf G_k)^T\in\mathbb C^{1\times I_1I_2}$ in the same
column-major order. Reshaping it back gives
\begin{equation}
\operatorname{mat}\!\left([G_{(3)}]_{k,:}\right)
=\mathbf G_k
\in\mathbb C^{I_1\times I_2},
\tag{402}
\end{equation}
where $\operatorname{mat}(\cdot)$ is the inverse of $\operatorname{vec}(\cdot)$
(in MATLAB, \texttt{reshape(row, I1, I2)}). Hence, from (382),
\begin{equation}
\operatorname{rank}\!\left(
\operatorname{mat}\!\left([G_{(3)}]_{k,:}\right)
\right)\leq L_k.
\tag{403}
\end{equation}

The channel estimation problem is therefore the rank-constrained
least-squares problem
\begin{equation}
\boxed{
\begin{aligned}
\min_{G_{(3)}}\quad
&f(G_{(3)})
=\left\|
Y_{(3)}-|B_{(3)}|
-\operatorname{Re}\!\left[
Z_{(3)}\circ(S^TG_{(3)})
\right]
\right\|_F^2\\
\text{s.t.}\quad
&\operatorname{rank}\!\left(
\operatorname{mat}([G_{(3)}]_{k,:})
\right)\leq L_k,
\qquad k=1,\ldots,K,
\end{aligned}
}
\tag{404}
\end{equation}
which is Eq.~(21) of Xu \emph{et al.} The rank constraint is what
distinguishes the 2D problem from the 1D problem (173), and it is why
the GS algorithm of Section~\ref{sec:gs} cannot be applied
directly.

% ============================================================
\section{Projected Gradient Descent for the 2D Channel}
% ============================================================

Define the residual
\begin{equation}
\boxed{
E
=
Y_{(3)}-|B_{(3)}|
-\operatorname{Re}\!\left[
Z_{(3)}\circ(S^TG_{(3)})
\right]
\in\mathbb R^{P\times I_1I_2}.
}
\tag{405}
\end{equation}

Then the objective in (404) is
\begin{equation}
f(G_{(3)})
=\|E\|_F^2
=\sum_{p=1}^{P}\sum_{q=1}^{I_1I_2}E_{p,q}^2,
\tag{406}
\end{equation}
which is Eq.~(22) of Xu \emph{et al.} Because $Y_{(3)}$, $|B_{(3)}|$,
and the real-part term are real, $E_{p,q}\in\mathbb R$.

\subsection{Elementwise Gradient Derivation}

Fix a differentiation index $(k,q)$ and use $r$ as the dummy user
index. The $(p,q)$-th element of $S^TG_{(3)}$ is
\begin{equation}
(S^TG_{(3)})_{p,q}
=
\sum_{r=1}^{K}[S^T]_{p,r}[G_{(3)}]_{r,q}
=
\sum_{r=1}^{K}s_{r,p}G_{r,q}.
\tag{407}
\end{equation}

Using $\operatorname{Re}(x)=(x+x^*)/2$,
\begin{equation}
\operatorname{Re}\!\left[
Z_{p,q}(S^TG_{(3)})_{p,q}
\right]
=
\frac12
\left[
Z_{p,q}\sum_{r=1}^{K}s_{r,p}G_{r,q}
+Z_{p,q}^{*}\sum_{r=1}^{K}s_{r,p}^{*}G_{r,q}^{*}
\right].
\tag{408}
\end{equation}

In Wirtinger calculus, $G_{r,q}$ and $G_{r,q}^{*}$ are treated as
independent variables, so that
\[
\frac{\partial G_{r,q}}{\partial G_{k,q}^{*}}=0,
\qquad
\frac{\partial G_{r,q}^{*}}{\partial G_{k,q}^{*}}=\delta_{r,k}.
\]
Only the second sum in (408) contributes, and only its $r=k$ term
survives. Therefore,
\begin{equation}
\boxed{
\frac{\partial E_{p,q}}
{\partial G_{k,q}^{*}}
=
-\frac12 Z_{p,q}^{*}s_{k,p}^{*}.
}
\tag{409}
\end{equation}

Since
\begin{equation}
f=\sum_{p,q}E_{p,q}^{2},
\tag{410}
\end{equation}
and only the entries in column $q$ of $E$ depend on $G_{k,q}$, the
chain rule gives
\begin{equation}
\begin{aligned}
\frac{\partial f}
{\partial G_{k,q}^{*}}
&=
2\sum_{p=1}^{P}
E_{p,q}
\frac{\partial E_{p,q}}
{\partial G_{k,q}^{*}}
=
2\sum_{p=1}^{P}
E_{p,q}
\left(-\frac12 Z_{p,q}^{*}s_{k,p}^{*}\right)
=
-\sum_{p=1}^{P}
s_{k,p}^{*}E_{p,q}Z_{p,q}^{*}.
\end{aligned}
\tag{411}
\end{equation}

Since $[S^{*}]_{k,p}=s_{k,p}^{*}$ and $[E\circ Z_{(3)}^{*}]_{p,q}=E_{p,q}Z_{p,q}^{*}$,
the sum in (411) is the $(k,q)$-th element of the matrix product
$S^{*}(E\circ Z_{(3)}^{*})$. Collecting all $(k,q)$,
\begin{equation}
\boxed{
\nabla_{G_{(3)}^{*}}f
=
-S^{*}\left(E\circ Z_{(3)}^{*}\right).
}
\tag{412}
\end{equation}

The dimensions are
\begin{equation}
\underbrace{S^{*}}_{K\times P}
\underbrace{\left(E\circ Z_{(3)}^{*}\right)}_{P\times I_1I_2},
\qquad
E\in\mathbb R^{P\times I_1I_2},
\quad
Z_{(3)}^{*}\in\mathbb C^{P\times I_1I_2},
\tag{413}
\end{equation}
so that
\begin{equation}
\nabla_{G_{(3)}^{*}}f
\in\mathbb C^{K\times I_1I_2},
\tag{414}
\end{equation}
which is the same dimension as $G_{(3)}$. Equation (412) is Eq.~(24) of
Xu \emph{et al.}, and it is the transposed 2D counterpart of the 1D
gradient (206). Like (206), it uses the Wirtinger convention; Eq.~(15)
of Xu \emph{et al.} carries an extra factor of $2$, which only rescales
the step size.

\subsection{Gradient Descent Step}

At iteration $t$, let
\begin{equation}
E^{(t)}
=
Y_{(3)}-|B_{(3)}|
-\operatorname{Re}\!\left[
Z_{(3)}\circ(S^TG_{(3)}^{(t)})
\right]
\in\mathbb R^{P\times I_1I_2}.
\tag{415}
\end{equation}

The gradient is
\begin{equation}
\boxed{
\nabla^{(t)}
=
-S^{*}\left(E^{(t)}\circ Z_{(3)}^{*}\right)
\in\mathbb C^{K\times I_1I_2}.
}
\tag{416}
\end{equation}

The unconstrained gradient step (step~4 of Algorithm~1 of Xu
\emph{et al.}) is
\begin{equation}
\boxed{
M
=
G_{(3)}^{(t)}
-\zeta_t\nabla^{(t)}
}
\tag{417}
\end{equation}
and hence
\begin{equation}
\boxed{
M
=
G_{(3)}^{(t)}
+\zeta_t S^{*}
\left(E^{(t)}\circ Z_{(3)}^{*}\right)
\in\mathbb C^{K\times I_1I_2},
}
\tag{418}
\end{equation}
where $\zeta_t>0$ is the step size.

\textbf{Step size.} For any perturbation
$\Delta\in\mathbb C^{K\times I_1I_2}$, since $|Z_{p,q}|=1$ and
$|\operatorname{Re}(x)|\leq|x|$,
\[
\left\|\operatorname{Re}\!\left[Z_{(3)}\circ(S^T\Delta)\right]\right\|_F^2
\leq
\|S^T\Delta\|_F^2
\leq
\lambda_{\max}(S^*S^T)\,\|\Delta\|_F^2
=
\lambda_{\max}(SS^H)\,\|\Delta\|_F^2 .
\]
Hence the curvature of $f$ is bounded by $\lambda_{\max}(SS^H)$, and a
fixed step size
\begin{equation}
\boxed{
\zeta_t\leq\frac{1}{\lambda_{\max}(SS^H)}
}
\tag{418a}
\end{equation}
is a sufficient (conservative) condition for a non-increasing objective in the
unconstrained step under the Wirtinger convention of (412).

The gradient step does not by itself enforce the rank constraints;
$\operatorname{mat}([M]_{k,:})$ is in general full rank.

\subsection{Projection onto the Rank-Constrained Set}

Define the feasible set
\begin{equation}
\boxed{
\mathcal X
=
\left\{
X\in\mathbb C^{K\times I_1I_2}:
\operatorname{rank}\!\left(
\operatorname{mat}([X]_{k,:})
\right)\leq L_k,
\ \forall k
\right\}.
}
\tag{419}
\end{equation}

The projected gradient update, Eq.~(23) of Xu \emph{et al.}, is
\begin{equation}
\boxed{
G_{(3)}^{(t+1)}
=
\operatorname{proj}_{\mathcal X}(M),
}
\tag{420}
\end{equation}
where, as in Eq.~(25) of Xu \emph{et al.},
\begin{equation}
\boxed{
\operatorname{proj}_{\mathcal X}(M)
=
\arg\min_{X\in\mathcal X}
\|X-M\|_F^2.
}
\tag{421}
\end{equation}

The Frobenius norm separates over the rows:
\begin{equation}
\|X-M\|_F^2
=
\sum_{k=1}^{K}
\|[X]_{k,:}-[M]_{k,:}\|_2^2.
\tag{422}
\end{equation}

Since the $k$-th constraint in (419) involves only the $k$-th row,
each term of (422) can be minimized independently, and each user's rank
constraint is enforced separately.

For user $k$, reshape the row using the same column-major ordering as
the mode-3 unfolding:
\begin{equation}
\boxed{
M_k
=
\operatorname{mat}([M]_{k,:})
\in\mathbb C^{I_1\times I_2}.
}
\tag{423}
\end{equation}

Since $\|[X]_{k,:}-[M]_{k,:}\|_2=\|X_k-M_k\|_F$ with
$X_k=\operatorname{mat}([X]_{k,:})$, the $k$-th subproblem is
\begin{equation}
\min_{\operatorname{rank}(X_k)\leq L_k}
\|X_k-M_k\|_F^2 .
\tag{423a}
\end{equation}

Compute the singular value decomposition, Eq.~(26) of Xu \emph{et al.}:
\begin{equation}
\boxed{
M_k
=
U_k\Sigma_kV_k^{H},
\qquad
U_k\in\mathbb C^{I_1\times I_1},
\ \
\Sigma_k\in\mathbb R^{I_1\times I_2},
\ \
V_k\in\mathbb C^{I_2\times I_2}.
}
\tag{424}
\end{equation}

Let the singular values be ordered as
\begin{equation}
\Sigma_k
=\operatorname{diag}(\sigma_{k,1},\sigma_{k,2},\ldots),
\qquad
\sigma_{k,1}\geq\sigma_{k,2}\geq\cdots\geq0.
\tag{425}
\end{equation}

Retaining the largest $L_k$ singular values gives, as in Eq.~(27) of
Xu \emph{et al.},
\begin{equation}
\boxed{
M_k^{(L_k)}
=
\underbrace{U_k(:,1:L_k)}_{I_1\times L_k}\,
\underbrace{\Sigma_k(1:L_k,1:L_k)}_{L_k\times L_k}\,
\underbrace{V_k(:,1:L_k)^H}_{L_k\times I_2}
\in\mathbb C^{I_1\times I_2}.
}
\tag{426}
\end{equation}

By the Eckart--Young--Mirsky theorem, $M_k^{(L_k)}$ is the closest
matrix to $M_k$ in the Frobenius norm among matrices of rank at most
$L_k$, i.e.\ it solves (423a):
\begin{equation}
\boxed{
M_k^{(L_k)}
=
\arg\min_{\operatorname{rank}(X_k)\leq L_k}
\|X_k-M_k\|_F^2,
\qquad
\|M_k^{(L_k)}-M_k\|_F^2=\sum_{r>L_k}\sigma_{k,r}^2.
}
\tag{427}
\end{equation}

Since $V_k$ is complex, the Hermitian transpose $V_k^H$ is used in
(424) and (426); Eqs.~(26)--(27) of Xu \emph{et al.} write $V_k^T$. If
$L_k\geq\min(I_1,I_2)$, the projection leaves $M_k$ unchanged.

After the projection, the $k$-th row of the updated unfolded channel is
\begin{equation}
\boxed{
[G_{(3)}^{(t+1)}]_{k,:}
=
\operatorname{vec}(M_k^{(L_k)})^T
\in\mathbb C^{1\times I_1I_2}.
}
\tag{428}
\end{equation}

Thus, one PGD iteration is
\begin{equation}
\boxed{
G_{(3)}^{(t)}
\xrightarrow{\text{gradient step (418)}}
M
\xrightarrow{\text{rank projection (426)--(428)}}
G_{(3)}^{(t+1)}.
}
\tag{429}
\end{equation}

\subsection{Complete PGD Algorithm}

The iterations stop when $\|G_{(3)}^{(t+1)}-G_{(3)}^{(t)}\|_F^2<\epsilon$
or $t\geq T_{\max}$. Following Xu \emph{et al.}, the initial estimate has
i.i.d.\ entries $[G_{(3)}^{(0)}]_{k,q}\sim\mathcal{CN}(0,0.1)$. The
complete procedure (Algorithm~1 of Xu \emph{et al.}) is
\begin{equation}
\fbox{%
\begin{minipage}{0.8\linewidth}
\vspace{4pt}
\textbf{Input:} $Y_{(3)},\ Z_{(3)},\ B_{(3)},\ S,\ \{L_k\}_{k=1}^{K},\ \zeta_t,\ \epsilon,\ T_{\max}$\\[4pt]
\textbf{Initialize:} $G_{(3)}^{(0)}$, \ $t=0$\\[4pt]
\textbf{Repeat:}\\[2pt]
\hspace*{1.5em}$E^{(t)}=Y_{(3)}-|B_{(3)}|-\operatorname{Re}[Z_{(3)}\circ(S^TG_{(3)}^{(t)})]$\\[2pt]
\hspace*{1.5em}$M=G_{(3)}^{(t)}+\zeta_t S^{*}(E^{(t)}\circ Z_{(3)}^{*})$\\[2pt]
\hspace*{1.5em}For $k=1,\ldots,K$:\\[2pt]
\hspace*{3em}$M_k=\operatorname{mat}([M]_{k,:})$, \ $M_k=U_k\Sigma_kV_k^H$\\[2pt]
\hspace*{3em}$[G_{(3)}^{(t+1)}]_{k,:}=\operatorname{vec}\!\big(U_k(:,1{:}L_k)\Sigma_k(1{:}L_k,1{:}L_k)V_k(:,1{:}L_k)^H\big)^T$\\[2pt]
\hspace*{1.5em}End\\[2pt]
\hspace*{1.5em}$t\leftarrow t+1$\\[4pt]
\textbf{Until:} $\|G_{(3)}^{(t)}-G_{(3)}^{(t-1)}\|_F^2<\epsilon$ \ or \ $t\geq T_{\max}$\\[4pt]
\textbf{Output:} $\widehat G_{(3)}=G_{(3)}^{(t)}$.
\vspace{4pt}
\end{minipage}}
\tag{429a}
\end{equation}

Per iteration, the gradient costs $\mathcal O(PKI_1I_2)$ and the $K$
SVDs cost $\mathcal O(K\min(I_1I_2^2,I_1^2I_2))$, which is the complexity
stated by Xu \emph{et al.} Without the projection, (429a) reduces to
plain gradient descent (GD).

% ============================================================
\section{CRLB Analysis for the 2D Channel}
% ============================================================

The Cram\'er--Rao lower bound is used as an estimation-accuracy
benchmark. The derivation below first follows the convention used by
Xu \emph{et al.} and then gives the exact real-parameter Fisher
information for the real observation model.

\subsection{Vectorized Observation Model}

Let
\begin{equation}
N_a=I_1I_2
\tag{430}
\end{equation}
denote the total number of array elements. To obtain a Kronecker
representation consistent with the mode-3 unfolding, define
\begin{equation}
\boxed{
 y=\operatorname{vec}(Y_{(3)})
 \in\mathbb R^{PN_a\times1},
\qquad
 g=\operatorname{vec}(G_{(3)})
 \in\mathbb C^{KN_a\times1}.
}
\tag{431}
\end{equation}

Here $\operatorname{vec}(\cdot)$ stacks the columns, so $y$ lists the
$P$ pilot measurements of antenna 1, then of antenna 2, and so on, and
$g$ lists the $K$ user channels of antenna 1, then of antenna 2, and so
on. Similarly,
\begin{equation}
 \mathbf z=\operatorname{vec}(Z_{(3)}),
\quad
 b=\operatorname{vec}(B_{(3)}),
\quad
 n_r=\operatorname{vec}(N_{(3)}),
\qquad
\mathbf z,b\in\mathbb C^{PN_a\times1},
\quad
n_r\in\mathbb R^{PN_a\times1},
\tag{432}
\end{equation}
where each entry of $\mathbf z$ has the form $e^{-jz_{i_1,i_2,p}}$ and
therefore $|[\mathbf z]_m|=1$.

Using the identity
\begin{equation}
\operatorname{vec}(AXB)
=
(B^T\otimes A)\operatorname{vec}(X),
\tag{433}
\end{equation}
with $A=S^T$, $X=G_{(3)}$ and $B=I_{N_a}$,
\begin{equation}
\boxed{
\operatorname{vec}(S^TG_{(3)})
=
(I_{N_a}\otimes S^T)g.
}
\tag{434}
\end{equation}

Using $\mathbf z\circ x=\operatorname{diag}(\mathbf z)x$, define
\begin{equation}
\boxed{
 A
=
\operatorname{diag}(\mathbf z)
(I_{N_a}\otimes S^T).
}
\tag{435}
\end{equation}

The dimensions are
\begin{equation}
\operatorname{diag}(\mathbf z)\in\mathbb C^{PN_a\times PN_a},
\qquad
I_{N_a}\otimes S^T
\in\mathbb C^{PN_a\times KN_a},
\qquad
A\in\mathbb C^{PN_a\times KN_a}.
\tag{436}
\end{equation}

Vectorizing (401), the observation model becomes
\begin{equation}
\boxed{
 y
 =
\operatorname{Re}(Ag)+|b|+n_r,
}
\tag{437}
\end{equation}
which is Eq.~(28) of Xu \emph{et al.}

From (391), the effective real noise satisfies
\begin{equation}
\boxed{
 n_r\sim\mathcal N(0,\sigma_r^2 I_{PN_a}),
 \qquad
 \sigma_r^2\triangleq\frac{\sigma^2}{2},
}
\tag{438}
\end{equation}
where $\sigma^2$ is the variance of the complex noise in (390), which
defines the SNR, and $\sigma_r^2$ is the variance of the effective real
noise.

\subsection{Likelihood Function}

The conditional mean of the observation is
\begin{equation}
\boxed{
\boldsymbol\mu(g)
=
\operatorname{Re}(Ag)+|b|
\in\mathbb R^{PN_a\times1}.
}
\tag{439}
\end{equation}

Since the entries of $n_r$ are independent Gaussian, the likelihood,
Eq.~(29) of Xu \emph{et al.}, is
\begin{equation}
\boxed{
 p(y|g)
 =
 \frac{1}
 {(2\pi\sigma_r^2)^{PN_a/2}}
 \exp\!\left[
 -\frac{1}{2\sigma_r^2}
 \|y-\boldsymbol\mu(g)\|_2^2
 \right].
}
\tag{440}
\end{equation}

Let
\begin{equation}
 e=y-\boldsymbol\mu(g)\in\mathbb R^{PN_a\times1}.
\tag{441}
\end{equation}

Then
\begin{equation}
\boxed{
\ln p(y|g)
=C-
\frac{1}{2\sigma_r^2}e^Te,
}
\tag{442}
\end{equation}
where $C=-\frac{PN_a}{2}\ln(2\pi\sigma_r^2)$ is independent of $g$.

\subsection{Xu-Compatible Fisher Information and CRLB}

Since
\begin{equation}
\operatorname{Re}(Ag)
=
\frac12(Ag+A^*g^*),
\tag{443}
\end{equation}
the Wirtinger derivatives of the mean are
\begin{equation}
\boxed{
\frac{\partial\boldsymbol\mu}{\partial g^T}
=\frac12A,
\qquad
\frac{\partial\boldsymbol\mu}{\partial g^H}
=\frac12A^*.
}
\tag{444}
\end{equation}

Since $\partial e/\partial g^H=-\tfrac12A^*$ and $e$ is real, the score
with respect to $g^*$ is
\begin{equation}
\boxed{
\frac{\partial\ln p(y|g)}{\partial g^*}
=
-\frac{1}{2\sigma_r^2}\,2\left(\frac{\partial e}{\partial g^H}\right)^{T}e
=
\frac{1}{2\sigma_r^2}A^He
\in\mathbb C^{KN_a\times1}.
}
\tag{445}
\end{equation}

Since $e=n_r$ at the true $g$,
\begin{equation}
\mathbb E[ee^T]
=\sigma_r^2I_{PN_a}.
\tag{446}
\end{equation}

The Fisher information block used by Xu \emph{et al.} is therefore
\begin{equation}
\begin{aligned}
F_{gg^*}
&=
\mathbb E\!\left[
\frac{\partial\ln p}{\partial g^*}
\left(
\frac{\partial\ln p}{\partial g^*}
\right)^H
\right]
=
\frac{1}{4\sigma_r^4}
A^H
\mathbb E[ee^T]
A
=
\frac{1}{4\sigma_r^2}A^HA.
\end{aligned}
\tag{447}
\end{equation}

Substituting (435),
\begin{equation}
A^HA
=
(I_{N_a}\otimes S^T)^H
\operatorname{diag}(\mathbf z)^H
\operatorname{diag}(\mathbf z)
(I_{N_a}\otimes S^T).
\tag{448}
\end{equation}

Since $|[\mathbf z]_m|=1$ for every entry,
\begin{equation}
\operatorname{diag}(\mathbf z)^H\operatorname{diag}(\mathbf z)
=\operatorname{diag}\!\left(|[\mathbf z]_1|^2,\ldots,|[\mathbf z]_{PN_a}|^2\right)
=I_{PN_a}.
\tag{449}
\end{equation}

Therefore, using $(I\otimes S^T)^H=I\otimes S^*$ and the mixed-product
property $(A\otimes B)(C\otimes D)=(AC)\otimes(BD)$,
\begin{equation}
A^HA
=
(I_{N_a}\otimes S^*)
(I_{N_a}\otimes S^T)
=
I_{N_a}\otimes(S^*S^T),
\qquad
S^*S^T\in\mathbb C^{K\times K},
\tag{450}
\end{equation}
and
\begin{equation}
\boxed{
F_{gg^*}
=
\frac{1}{4\sigma_r^2}
\left[
I_{N_a}\otimes(S^*S^T)
\right]
\in\mathbb C^{KN_a\times KN_a},
}
\tag{451}
\end{equation}
which is Eq.~(30) of Xu \emph{et al.}

Using $(A\otimes B)^{-1}=A^{-1}\otimes B^{-1}$, the inverse exists when
$S$ has full row rank ($P\geq K$), and
\begin{equation}
\operatorname{CRB}_{\mathrm{Xu}}
=F_{gg^*}^{-1}
=
4\sigma_r^2
\left[
I_{N_a}\otimes(S^*S^T)^{-1}
\right].
\tag{452}
\end{equation}

Hence the Xu-compatible CRLB is
\begin{equation}
\boxed{
\operatorname{CRB}_{\mathrm{Xu}}
=
4\sigma_r^2
\left[
I_{I_1I_2}\otimes(S^*S^T)^{-1}
\right],
}
\tag{453}
\end{equation}
which is Eq.~(31) of Xu \emph{et al.}, whose Eq.~(28) denotes the
effective real noise variance by $\sigma^2$ (i.e.\ their $\sigma^2$ is
$\sigma_r^2$ here). Expressed in the complex noise variance of (390),
which defines the SNR,
\begin{equation}
\boxed{
\operatorname{CRB}_{\mathrm{Xu}}
=
2\sigma^2
\left[
I_{I_1I_2}\otimes(S^*S^T)^{-1}
\right].
}
\tag{454}
\end{equation}

For example, if $I_1=I_2=8$ and $K=3$, then
\begin{equation}
N_a=64,
\qquad
KN_a=192,
\qquad
g\in\mathbb C^{192\times1},
\tag{455}
\end{equation}
and therefore
\begin{equation}
F_{gg^*},\operatorname{CRB}_{\mathrm{Xu}}
\in\mathbb C^{192\times192}.
\tag{456}
\end{equation}

The mean-square-error benchmark is the trace of the CRLB. Since
$\operatorname{tr}(I_{N_a}\otimes X)=N_a\operatorname{tr}(X)$,
\begin{equation}
\boxed{
\mathrm{MSE}_{\mathrm{CRB}}
=
\operatorname{tr}\!\left(\operatorname{CRB}_{\mathrm{Xu}}\right)
=
2\sigma^2 I_1I_2\,
\operatorname{tr}\!\left[(S^*S^T)^{-1}\right].
}
\tag{457}
\end{equation}
Because the mean in (437) is not holomorphic in $g$, (457) is the Xu-compatible
benchmark; the rigorous inequality for the real observation model is (463a).

The NMSE of Xu \emph{et al.} is
$\mathbb E\|\hat G_{(3)}-G_{(3)}\|_F^2/\mathbb E\|G_{(3)}\|_F^2$. Since
$\|G_{(3)}\|_F^2=\|g\|_2^2$, the CRLB benchmark curve is obtained by
normalizing (457) by the expected channel energy:
\begin{equation}
\boxed{
\mathrm{NMSE}_{\mathrm{CRB}}
=
\frac{\operatorname{tr}\!\left(\operatorname{CRB}_{\mathrm{Xu}}\right)}
{\mathbb E\|g\|_2^2}.
}
\tag{457a}
\end{equation}

Under the assumed channel statistics (equal average channel power per array
element), both the numerator and $\mathbb E\|g\|_2^2$ grow linearly with
$I_1I_2$, so $\mathrm{NMSE}_{\mathrm{CRB}}$ does not depend on the
number of antennas; it decreases with SNR through $\sigma^2$ and with
the pilot length through $\operatorname{tr}[(S^*S^T)^{-1}]$, as stated
by Xu \emph{et al.}

\subsection{Exact Real-Parameter Fisher Information}

The real-part operation in (437) makes the mean non-holomorphic in the
complex parameter $g$. Therefore, (453)--(454) are the Xu-compatible
Fisher-information block rather than an exact complex-parameter CRLB.
The following exact real-parameter formulation goes beyond the
derivation of Xu \emph{et al.} Write
\begin{equation}
\boxed{
 g=g_R+jg_I,
 \qquad
 A=A_R+jA_I,
}
\tag{458}
\end{equation}
with $g_R,g_I\in\mathbb R^{KN_a\times1}$ and
$A_R,A_I\in\mathbb R^{PN_a\times KN_a}$. Then
\begin{equation}
\operatorname{Re}(Ag)
=
\operatorname{Re}\!\left[(A_R+jA_I)(g_R+jg_I)\right]
=
A_Rg_R-A_Ig_I.
\tag{459}
\end{equation}

Define the real parameter vector and real Jacobian
\begin{equation}
\boxed{
\vartheta
=
\begin{bmatrix}
 g_R\\g_I
\end{bmatrix}
\in\mathbb R^{2KN_a\times1},
\qquad
D
=
\begin{bmatrix}
 A_R & -A_I
\end{bmatrix}
\in\mathbb R^{PN_a\times 2KN_a}.
}
\tag{460}
\end{equation}

The mean can then be written as a real linear model,
\begin{equation}
\boxed{
\boldsymbol\mu(\vartheta)
=D\vartheta+|b|.
}
\tag{461}
\end{equation}

For a real Gaussian linear model with noise covariance
$\sigma_r^2I$, the exact Fisher information matrix is
\begin{equation}
\boxed{
F_R
=
\frac{1}{\sigma_r^2}D^TD
\in\mathbb R^{2KN_a\times 2KN_a}.
}
\tag{462}
\end{equation}

If $F_R$ is nonsingular, the corresponding real-parameter CRLB is
\begin{equation}
\boxed{
\operatorname{CRB}_R
=F_R^{-1}
=\sigma_r^2\left(D^TD\right)^{-1}.
}
\tag{463}
\end{equation}

For $I_1=I_2=8$ and $K=3$, $\vartheta\in\mathbb R^{384\times1}$ and
$F_R,\operatorname{CRB}_R\in\mathbb R^{384\times384}$, compared with the
$192\times192$ complex matrices in (456). Since
$\|\hat g-g\|_2^2=\|\hat g_R-g_R\|_2^2+\|\hat g_I-g_I\|_2^2$,
\begin{equation}
\boxed{
\mathbb E\|\hat g-g\|_2^2
\geq
\operatorname{tr}\!\left(F_R^{-1}\right).
}
\tag{463a}
\end{equation}

Equivalently, with the augmented parameter $[g^T,\,g^H]^T$, the score is
$\frac{1}{2\sigma_r^2}[A^H;\,A^T]e$, and the augmented Fisher
information matrix is
\begin{equation}
\boxed{
\mathcal F_{\mathrm{aug}}
=
\frac{1}{4\sigma_r^2}
\begin{bmatrix}
A^HA & A^HA^*\\
A^TA & A^TA^*
\end{bmatrix}
\in\mathbb C^{2KN_a\times 2KN_a}.
}
\tag{464}
\end{equation}

The off-diagonal block is
\begin{equation}
A^HA^*
=
(I_{N_a}\otimes S^*)\,
\operatorname{diag}(\mathbf z^*\circ\mathbf z^*)\,
(I_{N_a}\otimes S^H),
\tag{464a}
\end{equation}
which is generally nonzero. For circularly symmetric pilots its expectation
is zero, and for a large pilot length $P$ it becomes negligible relative to
$A^HA$. When this block is zero,
$A_R^TA_R=A_I^TA_I=\tfrac12\operatorname{Re}(A^HA)$ and
$A_R^TA_I=-A_I^TA_R=\tfrac12\operatorname{Im}(A^HA)$, so that
\begin{equation}
\operatorname{tr}\!\left(F_R^{-1}\right)
=
\operatorname{tr}\!\left(\operatorname{CRB}_{\mathrm{Xu}}\right).
\tag{464b}
\end{equation}

When the off-diagonal block is nonzero, the $g$-block of
$\mathcal F_{\mathrm{aug}}^{-1}$ is the inverse of a Schur complement of
$F_{gg^*}$, which is larger than $F_{gg^*}^{-1}$. The Xu-compatible
bound (453) is therefore optimistic for short pilot sequences, while
(462)--(463a) give the exact bound for the real observation model.

\subsection{Relation Between PGD and the CRLB}

The PGD estimator incorporates the structural constraint
\begin{equation}
\operatorname{rank}(\mathbf G_k)\leq L_k,
\qquad k=1,\ldots,K,
\tag{465}
\end{equation}
whereas the CRLBs in (453) and (463) treat the unfolded channel
coefficients as unconstrained parameters. Therefore, the CRLB is an
accuracy benchmark for the unconstrained channel-estimation problem; a
rank-constrained CRLB would use $L_k(I_1+I_2-L_k)$ complex degrees of
freedom per user instead of $I_1I_2$. Both bounds apply to the
linearised model (437) and to unbiased estimators only; the model error
from the neglected second-order term in (387) is set by
$|(a_{i_1,i_2,p}+n_{i_1,i_2,p})/b_{i_1,i_2,p}|$, i.e.\ by the reference strength,
not by the SNR. It does not shrink as $\sigma^2\to0$, so at high SNR it
dominates the noise and appears as an NMSE floor for any estimator based on (437).

The number of real unknowns in an unconstrained complex channel vector
is $2KN_a$, and the number of real observations is $PN_a$. A generic
real-dimension counting condition for local identifiability of the
unconstrained channel is therefore
\begin{equation}
\boxed{
PN_a\geq2KN_a
\quad\Longrightarrow\quad
P\geq2K.
}
\tag{466}
\end{equation}

This is a generic counting condition rather than a sufficient condition
for every pilot/reference configuration. Identifiability also requires
$D$ in (460) to have full column rank.
In particular, for \(K\leq P<2K\), \(F_{gg^*}\) in (451) is invertible, so the
Xu-compatible bound (453) is finite, whereas \(F_R\) in (462) is singular and the
exact bound is infinite. This is a further reason why (453) is optimistic for
short pilot sequences.

% ============================================================
% END OF 2D, PGD, AND CRLB DERIVATION
% ============================================================

\end{document}
